% =====================================================================
%  Revision draft --- BeyondHulten / Metroeconomica
%  Working title: Which Short Run? Which Money? Closure Choice in
%                 Input-Output Models of the Socio-Ecological Transformation.
%  Basis: git/ai/RA/publishing/TeXTemplate (universal research-paper template)
%  Compile: latexmk -pdf main.tex   (see README.md for the container recipe)
% =====================================================================
\documentclass{article}

% ---- Switches (flip per venue) --------------------------------------
\newif\ifreviewmode   % true = double spacing + visible todonotes (submission)
\newif\ifendfloats    % true = floats deferred to end of document
\reviewmodetrue
\endfloatstrue

% ---- Page ------------------------------------------------------------
\usepackage[a4paper, total={6in, 9in}]{geometry}
\usepackage[T1]{fontenc}
\usepackage{parskip}
\usepackage{setspace}
\ifreviewmode\doublespacing\fi

% ---- Graphics & drawing ----------------------------------------------
\usepackage{graphicx}
\graphicspath{{figures/}}
\usepackage{tikz}
\usetikzlibrary{arrows.meta, positioning, calc}
\usepackage{subcaption}
\usepackage{float}
\usepackage{lscape}

% ---- Math & theorems --------------------------------------------------
\usepackage{amsmath, amssymb, amsthm}
\usepackage{booktabs}
\newtheorem{theorem}{Theorem}
\newtheorem{proposition}{Proposition}
\newtheorem{lemma}{Lemma}
\newtheorem{corollary}{Corollary}
\newtheorem{definition}{Definition}
\newtheorem{assumption}{Assumption}
\newtheorem{remark}{Remark}

% macros used by Appendix A (from equivalence.tex)
\newcommand{\Lbar}{\bar{L}}
\newcommand{\one}{\mathbf{1}}

% ---- Notes & boxes ----------------------------------------------------
\ifreviewmode
  \usepackage{todonotes}
\else
  \usepackage[disable]{todonotes}
\fi
\usepackage{framed}

% ---- Bibliography (biblatex + biber, authoryear) -----------------------
\usepackage[
backend=biber,
style=authoryear,
isbn=false, url=false, doi=false, eprint=false,
natbib=true,
giveninits=true,
uniquelist=false
]{biblatex}
\addbibresource{references.bib}

% ---- Authors & affiliations -------------------------------------------
\usepackage{authblk}

% ---- Floats to end of document ----------------------------------------
\ifendfloats
  \usepackage[markers,nolists]{endfloat}
\fi

% ---- Links & appendix --------------------------------------------------
\usepackage{hyperref}
\usepackage[toc,title]{appendix}

% ---- Date format -------------------------------------------------------
\renewcommand{\today}{~%
  \ifcase \month
    \or January\or February\or March\or April\or May\or June\or July\or August\or September\or October\or November\or December\fi,\space\number \year}

% ---- Title -------------------------------------------------------------
% TITLE SUGGESTIONS (from the dissolved old draft / the outline):
%  (1) working title (paper/narrative_outline.md): Which Short Run? Which Money? ...
%  (2) devised for the old draft: Labor Supply Elasticities and Sectoral Dynamics
%      in the Input-Output Modeling Continuum: A Quantitative Decomposition ...
\newcommand{\papertitle}{Which Short Run? Whose Money? Closure Choice in Input--Output Models of the Socio-Ecological Transformation}
\title{\papertitle\thanks{This paper has benefited from funding provided by the \textit{Hans B\"ockler Foundation} under grant number 2021-544-2, by the \textit{Linz Institute for Transformational Change (LIFT\_C)} under grant number 2022-1-114 as well as by \textit{Dezernat Zukunft}. Moreover, we thank Lukas Cserjan and Bernhard Sch\"utz for their extensive feedback as well as David Baqaee, Anna Hornykewycz, Stefan Steinerberger, Rafael Wildauer and Jan David Weber for their responses to our queries and the comments on specific aspects of this manuscript. All remaining errors are our own.}}

\author[1]{Jakob Kapeller}
\author[2]{Franz Scharnreitner}
\affil[1]{\small{Corresponding author, Institute for Socio-Economics, University of Duisburg-Essen and \linebreak Institute for Comprehensive Analysis of the Economy (ICAE), Johannes Kepler University Linz, \href{mailto:jakob.kapeller@uni-due.de}{jakob.kapeller@uni-due.de}}}
\affil[2]{\small{Institute for Comprehensive Analysis of the Economy (ICAE), Johannes Kepler University Linz}}
\date{\today}

\begin{document}
\maketitle

\begin{abstract}
% TODO: one paragraph. Headline = the transformation use case first; the
% equivalences as supporting rigour. Do NOT say "collapse theorem".
% Numbers per ADR-0024: G0 = 40.3 bn (2019 prices) = the horizon mean,
% 1.331 % of GDP; 58.9 bn/yr at 2023 prices; 2024 peak 81 bn ~ 1.9 % of GDP.
[Abstract to be drafted.]

Input-output models are a widely used tool for estimating the economic effects associated with societies moving towards carbon-neutrality or, more fundamentally, aiming to achieve a socio-ecological transformation. While this practice documents that such models are perceived as conceptually useful and empirically relevant, such approaches often take input-output models as a supposedly neutral analytical device. Against this background, this paper revisits how paradigmatic perspectives matter for input-output modeling. It shows that such perspectives and the associated analytical differences have a substantial impact on estimated effects at various levels. This observation indicates how modern applications of established models are actually haunted by ghosts from the past that reflect long-standing debates, and this has impact on the interpretation of a great variety of published results. Using the case of an archetypal green transformation policy, the decarbonization of the German housing sector, we show how different axiomatic setups imprint on estimates of the economic effects of ecologically motivated transformation efforts. By doing so, we hope to provide applied researchers with better guidance for reflecting on how the diversity of available interpretations of input-output models impacts on their problem at hand.

\textbf{Keywords:} input--output modeling, general equilibrium, closure choice, green investment, socio-ecological transformation

\textbf{JEL-Codes:} C67, C68, D57, E24
\end{abstract}

\newpage

%%% SECTION INPUTS %%%

% =====================================================================
% Section 1 --- Introduction
% ---------------------------------------------------------------------
% Job: transformation hook; the hidden closure choice; contributions
%      (rigidity map, closure taxonomy, 5x3 on a published impulse,
%      equivalences listed but not headlined); reading-guide box.
% Mindmap: N01, N08-N10.
% Status: STUB --- narrative outline v3 (2026-09-20).
% ---------------------------------------------------------------------
% SURVIVING TEXT from the old draft (paper/main.tex sec. 1): the Kuhnian
% cleavages / axiomatic-variation vocabulary --- reuse largely as-is, but
% keep the Kuhnian paragraph to a minimum (R2.3: reframe via domain of
% applicability; the cleavage framing is the reviewer's main objection).
% =====================================================================
\section{Introduction}\label{sec:intro}

% TODO (drafting notes):
% - Open concrete: the transformation programme (Sec. 2 gives the vector).
% - Name the hidden choice: the "multiplier" everyone quotes assumes a short
%   run and a financing regime; closure is *the* axiomatic variation.
% - Contributions: (i) the rigidity map; (ii) the closure taxonomy;
%   (iii) the 5x3 matrix on a published impulse; (iv) the equivalences
%   (listed, not headlined).
% - State upfront that the IO-CGE bridge is known (Robinson 2006; Rose 1995);
%   no bridge-discovery claim (R2.1/R2.3).
% - Reading-guide box: beginners Sections 1, 2, 5; experts Sections 3, 4, 6.

% TODO (reading-guide box --- framed environment):
% \begin{framed}
% \textbf{How to read this paper.} Beginners: Sections 1, 2 and 5.
% Experts: Sections 3, 4 and 6. Proofs are in Appendix A.
% \end{framed}
% =====================================================================
% Section 2 --- From transformation research to input-output models
% ---------------------------------------------------------------------
% Job: why *plural* IO models are the bridge; the construction and price
%      basis of the empirically informed green-investment vector; light
%      mathematics, the beginner on-ramp.
% Mindmap: N08, N09, N10.
% Status: STUB --- narrative outline v3 (2026-09-20).
% =====================================================================
% NOTE (numbers belong here, not in the abstract) -- per ADR-0024:
%   G0 = 40.3 bn in 2019 prices = the horizon mean (2024-2050) = 1.331 % of GDP
%        (40.3 / 3027.8; equals the booked B_gov = 0.01330991);
%   58.9 bn/yr in 2023 prices = the raw impulse mean (the CSV's own basis);
%   2024 first-year peak 81 bn (2023 prices) ~ 1.9 % of GDP;
%   deflator 1.46 (construction-cost indices); 58.89 / 1.46 = 40.338 = 40.337 bn.
% Cite BOTH price bases; never mix them in one sentence (outline external
% review, "The GDP share of the programme").
\section{From transformation research to input--output models}\label{sec:io-bridge}

% TODO (drafting notes):
% - Open concrete: the npj Climate Action impulse (Hornykewycz et al. 2025),
%   ~58 bn/yr and 3.1 tn until 2050 in 2023 prices; the 2024 first-year peak
%   ~81 bn ~ 1.9 % of GDP; G0 = 40.3 bn in 2019 prices = the horizon mean,
%   1.331 % of GDP (ADR-0024: both price bases cited, never mixed).
% - Why sectoral IO models are the applied lingua franca of transformation
%   research (Creutzig et al. 2018; Miller & Blair; Leontief 1936, 1986).
% - The vector: `impulses.csv`, incidence psi, the deflator 1.46 and the
%   ADR-0024 price-basis verification (defer the full provenance to App. D).
% - The beginner on-ramp: keep the algebra light here; the formal kernel is
%   Section 4 and Appendix B.

% ===== CARRIED OVER FROM AN OLD DRAFT =============================
% [devised for an old draft; not in the validated submission]
% source: revised_manuscript/chapters/ch04_data.tex (commit 41144dd)
% =================================================================
% To provide a concrete and reliable empirical example for such investment requirements, we take the German residential housing sector as our main case of analysis. When doing so we base our analysis and discussion on a specific estimate of the costs associated with a transformation of the German housing sector towards carbon neutrality as supplied by \citet{Hornykewycz.2025}. This study estimates expected costs by proportionally scaling past renovations costs under the assumption that an increase in the renovation intensity takes place, that is sufficient to meet Germany's climate goals. Specifically, the study at hand finds that renovation efforts will have to be more than doubled, if Germany is set to comply with its overall goal to achieve carbon neutrality till 2050. Such a shift would amount to yearly costs of about 40.3 bn. € in 2019 prices\footnote{We assumed a deflator of 1.46 based on price indices for construction costs as specified in \citet{Hornykewycz.2025}. In 2023 prices the shock amounts to about 58 bn. €.}, which would represent a substantial shift in final demand. As the study also details how these costs should be mapped on different sectors, it provides sufficient information to construct a concrete and directly applicable specification for the sector-specific annual changes in final demand imposed by such a reorientation. This specification shocks final demand in seven sectors, where sectors are sorted by the relative increase associated with these shocks (as given in brackets): 
% Glass and glassware (+50.3\%), ceramic products and processed stones (+37,3 \%), specialized construction work (+27,8\%), rubber and plastic products (+13.4\%), chemicals and chemical products (+3.7\%), machinery (+1.2\%) and electrical equipment (+1.0\%).
% In terms of the data used, we apply the so estimated shock vector as proposed by \citet{Hornykewycz.2025} with official German input-output data for 2019, the most recent year available \citep[see][]{destatis_strukturdaten}.


% =====================================================================
% REVIEW COMMENTS on this draft (agent; block refreshed 2026-10-11)
% ---------------------------------------------------------------------
% What works: the split is clean -- labour-market closures carry the vision
% of the short run, financial closures carry the theory of money -- and the
% money paragraph is still the sharpest statement of the R2.4 answer in the
% draft. The paragraph added today is a real improvement: it shows that the
% short-run axis is not one-dimensional, by putting a flexibility refinement
% and a rigidity refinement side by side.
%
% Changes made today by the agent (all reversible):
%   * typos fixed: modificaiton, aggreate, heoplessly, "from other sector",
%     and the mangled quotes around "no-reallocation".
%   * the footnote now opens with "the elastic labour supply closure"
%     rather than with an internal id (see the naming note below).
%
% Open points on today's text:
%  1. The paragraph covers two refinements (elastic supply, then
%     no-reallocation) but not the Leontief corner. A transition is drafted
%     in the block below the paragraph; it also supplies the "variety
%     collapses at the parameter ends" line, properly scoped.
%  2. "these different options on closure choices" is redundant; suggest
%     "these closure options".
%  3. "further established variation [of] closure setups exist" -- number
%     agreement; suggest "further established variations ... exist".
%  4. The table's column names ("standard CGE", "standard Leontief") are
%     never introduced in the text. One sentence after the table does it;
%     drafted in item E of the block below.
%  5. US/UK spelling is mixed through the draft ("labor"/"labour",
%     "modeling"/"modelling"). Metroeconomica accepts either, but not both.
%  6. Still flagged: the money claim ("roughly emulate an endogenous concept
%     of money") -- the must-not-appear list wants "an accounting analogue of
%     external finance"; and the two "(AUSNAHMEN ZITIEREN??)" gaps.
%
% NAMING: the paper should not wear the internal ids (ALPHA/BETA/GAMMA/
% DELTA/BF). The proposed vocabulary is in item E below; it needs one word
% of approval before it is swept through the draft.
%
% See PROPOSED FOOTNOTE / NOTES / SOURCE SUGGESTIONS at the foot of this file.
% =====================================================================
% Section 3 --- Two primitive questions
% ---------------------------------------------------------------------
% Job: 3.1 vision of the short run (quantity adjustment vs. rationing vs.
%      clearing, mapped via Taylor closures, FIDELIO, the neo-Keynesian
%      rigidity catalogue). 3.2 theory of money (F1/F2/F3 as endogenous-
%      money vs. balanced-budget vs. external-funds answers; F3 claimed as
%      an *accounting* closure, not a monetary mechanism). A stated as the
%      operational summary of 3.1 and 3.2.
% Mindmap: N02B, N02C, N02A.
% Status: STUB --- narrative outline v3 (2026-09-20).
% =====================================================================
\section{Two primitive questions, a simplified matrix and some complications}\label{sec:two-questions}

As already indicated, key points of disjuncture in input-output modeling can be mapped onto classic cleavages separating different traditions of economic thought. Such traditions can often be characterized in terms of the basic analytical ideas -- or ``rivers'' in the history of economic thought (SCREPANIT/ZAMGANI with p.) -- they endorse. Where these basic ideas have a different emphasis or contradict outright, a paradigmatic cleavage is in sight \citep{Dobusch.2012}. Two such classic fields of tension in economic thinking also lie at the heart of different approaches to input-output modeling: the nature of the short run and the role of money and credit, as already captured by Table~\ref{tab:overview}.

When mapping the rough distinctions from our introductory table more specifically to the realm of input-output models, we find that different labor market closures are the key entry point for different understandings of the short run: if labor is tight or fully utilized, increased monetary demand will increase wages and prices as capacity constraints bind. Conversely, if free labor is available, additional expenditures might lead to an expansion of employment and, correspondingly, output. Both visions can be unified in one analytical apparatus (see ROSBINSON UND PTHER POSER SOURCES? or \ref{sec:kernel}), just differentiated by different closures, i.e., definitions of what is exogenous or endogenous, in the face of increasing monetary expenditures.

Financial closures on the other hand encode prior beliefs about the role of our monetary system and the nature of money. As input-output models are typically conceived as static, operating outside of time (see however AUSNAHMEN ZITIEREN??), there is no borrowing against the future -- additional expenditure cannot be externally financed in the conventional sense, but have to be repaid within the same period (see however AUSNAHMEN ZITIEREN??). Hence, while a traditional, simplified application of a Leontief inverse would suggest some external stimulus, exactly such an external force is at odds with the internal consistency of the model. In other words, such models implicitly build on an exogenous view on money and restrict the application of ideas on endogenous money creation as it lacks a corresponding model of the intertemporal role of debt.

Refinancing can then occur in three ways: by households -- the repurposing of consumption --, via the state -- due to increased taxation --, or by someone willing to take foreign debt to provide an external stimulus. The latter can thereby be used to roughly emulate an endogenous concept of money within an input-output framework without explicitly violating its internal consistency properties. Mapping these two dimensions onto specific modeling approaches allows for recasting Table~\ref{tab:overview} in terms of specific constellations of closure summarized by Table~\ref{tab:closure-conditions} below.

% =====================================================================
% Table 2 --- the closure conditions, in simplified formal form
% ---------------------------------------------------------------------
% What Section 3's closing sentence promises: Table 1 recast "in terms of
% specific constellations of closure". Here the content is the CLOSURE
% CONDITIONS themselves -- two dimensions, five closures -- with the equation
% each one pins and the margin that is left to adjust.
% Alternative shape (if you prefer the 3x2 grid kept): the same conditions in
% the row/column headers of Table 1, with the qualitative outcomes in the
% cells. Say the word and I swap it.
% =====================================================================
\begin{table}[htbp]
\centering
\small
\begin{tabular}{@{}l l l@{}}
\toprule
\textbf{Closure} & \textbf{Condition (simplified)} & \textbf{What adjusts} \\
\midrule
\multicolumn{3}{@{}l}{\textit{Short run --- what the labour market pins}} \\[2pt]
factor supply fixed
  & $\sum_i L_i = \bar L$
  & wages and prices \\[3pt]
real wage fixed
  & $w/P = \bar w$
  & employment $L$ \\[5pt]
\multicolumn{3}{@{}l}{\textit{Money --- how the programme is financed}} \\[2pt]
repurposing of consumption (F1)
  & $\sum_i p_i c_i = E_h$
  & the composition of $c$ \\[3pt]
refinancing via taxes (F2)
  & $\sum_i p_i g_i = T$
  & the tax bill $T$ \\[3pt]
foreign debt (F3)
  & $\sum_i p_i g_i = F$
  & the external transfer $F$ \\
\bottomrule
\end{tabular}
\caption{Table~\ref{tab:overview} recast as specific constellations of
closure: the conditions each closure pins, in simplified form. The full
matrix follows in Section~\ref{sec:application}.}
\label{tab:closure-conditions}
\end{table}


While the crosswise recombination of these different options on closure choices already gives rise to some complexity, the underlying choices are somewhat rough and -- in the case of the short-run -- simplistically dichotomic. Hence, it comes as no surprise that further established variation of closure setups exist that try to address this point. An obvious starting point from a mainstream perspective is to introduce greater variety in labor market outcomes by explicitly introducing a leisure-income trade-off, which leads to the \textit{elastic labor supply model}: Assuming that the substitution effect dominates, higher wages will then attract additional labor so that capacity constraints do not bite with full force.\footnote{%TODO:diese fußnote reetten?
This is the \emph{elastic labour supply} closure: the fixed labour endowment is replaced by a supply curve, $L_i^{s} = \bar L_i\,[\,(w_i/\Pi(p))/(w_i^{0}/\Pi^{0})\,]^{\eta_s}$, with one common elasticity $\eta_s$, deflated by the consumption price index and anchored at $w_i^{0}/\Pi^{0} = 1$. Two changes of functional form follow: labour supply becomes \emph{behavioural} rather than a quantity restriction or a price pin, and the single aggregate labour market becomes $N$ sectoral markets with $N$ wages. The parameter range is $\eta_s \in [0, \infty)$, and it moves along the short-run axis of Table~\ref{tab:closure-conditions} instead of adding a third pole: a vertical supply curve at $\eta_s = 0$ freezes the allocation and lets the $N$ wages absorb (the quantity-rigid corner; in the single-wage form it is the full-employment pole instead), while $\eta_s \to \infty$ flattens the curve into a pinned real wage (the price-pinned corner). Its empirical content is therefore a slope, not a corner: on the executed ladder the price response falls monotonically, $\max|p-1|$ dropping from $0.153$ to $0.037$ as $\eta_s$ runs from $0.25$ to $2$, with the headline panels at $\eta_s = 0.5$. Two limits should be kept in view: under demand-only shocks a single economy-wide wage keeps the supply curve out of the price block whatever $\eta_s$, so the elasticity is an assumed sensitivity parameter rather than an estimate --- identifying it needs the supply-side arm of Section~\ref{sec:beyond-demand-only}; and the implemented curve is a reduced form without an income effect, so the leisure--income reading describes the slope of a supply relation, not a microfoundation.}. While this modification necessarily brings greater responsiveness on the side of employment other variations in the mainstream literature assume forms of rigidity that even go beyond the conventional CGE labor market closure. A prominent recent example is given by \citet{BF2019}, who introduce a ``no-reallocation'' case as a baseline scenario to articulate the view that sectoral productivity shocks might impact on aggregate outcomes in a non-linear way. In the context of our case of application -- a socio-ecological transformation that is tied to sector-specific shifts or injections in final demand -- such a baseline setup is especially pessimistic as it assumes labor market rigidity across sectors, which are then hopelessly overwhelmed as workers from other sectors cannot ``reallocate'' to where capacity is needed.

% =====================================================================
% FRAMING: THE LEONTIEF CORNER, AND WHAT THE EQUIVALENCES ALREADY BUY US
% (agent comments, 2026-10-11 -- no text changed; adopt, edit, or drop)
% ---------------------------------------------------------------------
% A. WHERE THE PARAGRAPH STANDS
%    It moves from "the short run is dichotomous" through two mainstream
%    refinements pointing in opposite directions: (i) elastic labour supply,
%    more responsiveness at the wage margin; (ii) BF (2019)'s
%    no-reallocation case, more rigidity at the allocation margin. That pair
%    is valuable precisely because it shows the axis is not one-dimensional.
%    Missing: the corner most applied IO work actually uses.
%
% B. SUGGESTED TRANSITION TO THE LEONTIEF CORNER  [draft, adopt or drop]
%
%    The two refinements just described move the labour market in opposite
%    directions: one buys flexibility at the wage margin, the other removes
%    it at the allocation margin. A third variation, and the oldest of the
%    three, pulls a different lever: it removes flexibility from production
%    itself. Replace CES substitution by fixed technical coefficients and
%    peg the real wage, and the kernel collapses to the classical Leontief
%    price-quantity system, in which prices follow from the cost identity
%    alone and final demand determines output. This is the corner most
%    applied transformation studies work with, usually without saying so.
%    It is not a new mechanism but a limit of the fixed-wage case: take its
%    substitution elasticities to zero and nothing is left to adjust on the
%    price side. It is also the mirror image of the no-reallocation case --
%    there quantities are frozen and prices absorb, here prices are frozen
%    and quantities absorb demand. The two therefore close the short-run
%    axis of Table~\ref{tab:closure-conditions}, and every further
%    refinement travels between them.
%
%    Why this framing works: it (i) motivates a THIRD entry without letting
%    the chapter become a list, (ii) connects to the BF sentence already in
%    the text by symmetry (quantities vs. prices), and (iii) sets up the
%    closing joke in C, which needs both ends established first.
%
% C. THE "FUNNY STATEMENT" -- THE VARIETY COLLAPSES AT THE PARAMETER ENDS
%
%    The honest version is unusually strong here, so it is worth saying
%    plainly -- but it must be scoped, or it is a false claim.
%
%    What holds (demand-only shocks, the branch the numeraire selects):
%      * The elastic-supply refinement adds a parameter, not an economy:
%        with a SINGLE economy-wide wage the supply curve is anchored at
%        the baseline real wage, so the anchor binds and the curve returns
%        the baseline quantity whatever its elasticity. The extra
%        flexibility is inert, and the closure coincides with the standard
%        CGE case (Appendix~\ref{app:proofs}, Lemma~\ref{lem:price}).
%      * Its rigid endpoint is the no-reallocation case: at eta_s = 0 the
%        supply curve is vertical in every sector, the allocation is frozen
%        and the N wages absorb -- exactly the case named just above. The
%        nesting is proved analytically and measured at 0.00e+00.
%      * Its elastic endpoint is the fixed-wage case: as eta_s -> infinity
%        the curve flattens into a pinned real wage. That is a limit, not a
%        cell -- the direct formulation stiffens and eventually fails.
%      * The Leontief corner adds nothing on top of the fixed-wage case:
%        under demand-only shocks a fixed wage already pins p = 1 exactly,
%        so the two coincide.
%    Hence: over its whole range the new parameter moves the economy ALONG
%    the axis, and at both ends it lands on a case we already had.
%
%    Candidate closing lines (pick one, or none):
%      (i)   "The menu of closures turns out to be longer than the list of
%             economies."
%      (ii)  "Every refinement we have introduced has an endpoint that is a
%             cell of the table we started from."
%      (iii) "Like a return ticket: an extra parameter carries the economy
%             away from a pole and, at sufficiently extreme values, brings
%             it back to one."
%
%    THE SCOPE CLAUSE IS NOT OPTIONAL -- with the joke, say:
%      1. demand-only shocks;
%      2. the single-wage form only: with N sectoral wages the elastic
%         family DOES move prices, and that is exactly where this paper's
%         variety lives (Section~\ref{sec:beyond-demand-only});
%      3. "collapse" concerns the PRICE margin, not the labour market:
%         employment still adjusts in the fixed-wage and Leontief cases.
%    Without (2) the sentence is false; without (1) it is unlicensed.
%
% D. EQUIVALENCE INVENTORY -- WHAT CAN BE SAID HERE, AND WHAT WAITS
%    (all "exact" entries are demand-only; Section 6 owns the proofs)
%
%      pair                                scope                status
%      ------------------------------------------------------------------
%      elastic supply = standard CGE        any eta_s, 1 wage    exact
%      elastic supply (eta_s = 0) = the     every sector         exact
%        no-reallocation case
%      elastic supply (eta_s -> inf) =      limit                asymptotic,
%        fixed real wage                                          never attained
%      Leontief = fixed real wage           demand-only          exact
%      F2 = F3 in the flexible-wage rows    demand-only          exact
%      F2 =/= F3 in the fixed-wage rows     demand-only          measured
%      ------------------------------------------------------------------
%    Suggested use in THIS section: one forward-pointing sentence, in the
%    neighbourhood of the joke -- "as Appendix~\ref{app:proofs} shows, and
%    Section~\ref{sec:beyond-demand-only} exploits, several of these cells
%    turn out to describe the same economy; the financing column is where
%    the choice really bites." Everything else belongs in Section 6; do not
%    unfold the propositions here.
%
% E. NAMING: CONVENTIONAL VOCABULARY INSTEAD OF THE INTERNAL IDS
%    [needs one word of approval before a sweep through the draft]
%
%      internal id   proposed paper name          literature anchor
%      ------------------------------------------------------------------
%      ALPHA         standard CGE closure         Sen (1963); Taylor &
%                    (neoclassical closure)      Lysy (1979); Robinson
%                                                (2006); the most common
%                                                factor-market closure
%      BETA          elastic labour supply;       Roy (1951); Lagakos &
%                    sectoral form: Roy-type     Waugh (2013); Galle et
%                    sectoral labour supply      al. (2015) -- the supply
%                                                elasticity to a sector
%      GAMMA         fixed real wage closure     "fixed real wage" /
%                                                neo-Keynesian closure II
%      DELTA         Leontief closure            Leontief (1936, 1986);
%                    (fixed coefficients)        Robinson (2006)
%      BF            no-reallocation closure     BF (2019); the
%                    (specific factors)          specific-factors case
%      ------------------------------------------------------------------
%    Note the asymmetry: DELTA is Leontief PLUS the fixed real wage, so its
%    name should not promise more than the technology. App. A and App. E
%    keep the ids as technical shorthand where they are defined; only the
%    PROSE renames. Your own sentence today ("beyond the conventional CGE
%    labor market closure") already uses the vocabulary.
%
%    Column sentence for after the table (item 4 of the review block):
%
%    The two columns are the two standard short-run closures of the
%    literature. The first is the neoclassical CGE closure, in which a
%    flexible wage clears the labour market at full employment and only
%    relative prices move. The second is the Leontief closure, in which
%    fixed technical coefficients and a fixed real wage leave prices pinned
%    and output to absorb demand. Applied transformation studies are built
%    on one or the other, usually without saying which.
% =====================================================================



% TODO: quantity adjustment vs. rationing vs. clearing; Taylor (1990)
% structuralist closures; the FIDELIO model (Rocchi et al.) as the applied
% neo-Keynesian alternative; the rigidity catalogue; the no-go corollary
% boundary (nominal wage rules cannot manufacture real variety --- forward
% reference to Sec. 6 and App. A).


% TODO: F1 preference reallocation (endogenous-money intuition, no net
% injection), F2 tax-financed (balanced budget), F3 external/debt (external
% funds; claim as an *accounting* closure, not a monetary mechanism).
% Robinson (2006) central; Lavoie (2014), Blecker (2019).

% TODO (closing): state the adjustment closure A as the operational summary
% of B (short run) and C (money); flag that A is *measured* from Sec. 4 on.
% =====================================================================
% PROPOSED FOOTNOTE (for para. 1; accept, drop, or move)
% ---------------------------------------------------------------------
% \footnote{We focus on these two dimensions because they are the closure
% choices that move transformation-programme estimates most in our setting,
% while the production-technology dimension is deliberately held fixed: under
% demand-only shocks the price block is demand-invariant, so the substitution
% elasticities do not reach the aggregate result (Appendix~\ref{app:proofs},
% Lemma~\ref{lem:price}). The remaining divides of Section~\ref{sec:intro}
% --- policy implementation and technological change --- enter only as
% parameters of the programme itself.}
% =====================================================================
%
% =====================================================================
% NOTES on the three related questions
% ---------------------------------------------------------------------
% (2) Kaldor-Verdoorn. It is measured but it is NOT a closure: it is one of
%     the seven fixed-wage (GAMMA) variants (probe17: the capacity channel
%     delta = +0.5 vs the Kaldor-Verdoorn case delta = -0.5; a 3.6 pp
%     deflator range from the sign alone). So it belongs in the GAMMA menu
%     of Section 6 / Appendix C and in the "open doors" of Section 7 -- not
%     here. Substitution elasticities: agreed, no need to dwell -- they are
%     inert under demand-only shocks (same Lemma 1 argument as the footnote
%     above) and matter only for the supply arm and the sectoral detail.
% (3) Equivalence appendix. Yes -- one light, scoped forward reference. The
%     sentence "Both visions can be unified in one analytical apparatus" IS
%     Lemma 1 + Propositions 1-3; add "(Appendix~\ref{app:proofs})" there and
%     keep the scope word "demand-only" next to it, since the unification
%     fails for supply shocks (Section~\ref{sec:beyond-demand-only}). Do not
%     unfold the equivalences here -- Section 6 owns them.
% =====================================================================
%
% =====================================================================
% SOURCE SUGGESTIONS for the bracketed placeholders
% (keys without NEW are already in revision/references.bib and compile)
% ---------------------------------------------------------------------
% (a) "(SCREPANIT/ZAMGANI with p.)":
%     \citep{Screpanti2005}   [= Screpanti & Zamagni, An Outline of the
%     History of Economic Thought; give the page for the "rivers" image]
% (b) "(DOBUSCH+KAPELLER12)":
%     \citep{Dobusch.2012}    [check the author list of that entry -- the
%     bib key is Dobusch.2012]
% (c) "(see ROBINSON UND PTHER POSER SOURCES? ...)":
%     NEW Robinson.2006 = Robinson, S. (2006), "Macro models and multipliers:
%     Leontief, Stone, Keynes, and CGE models", in de Janvry & Kanbur (eds),
%     Poverty, Inequality and Development, Springer, 205-232. Already
%     compile-ready alternatives: \citep{Taylor1990, lavoie2014, Blecker2019}.
% (d) "(AUSNAHMEN ZITIEREN??)"  (x2) -- the dynamic-IO exceptions:
%     \citep{BlairMiller} for the standard static treatment;
%     NEW candidates: Rocchi et al. (JRC FIDELIO manual) for the dynamic
%     macroeconometric IO alternative; Duchin (1998), Structural Economics,
%     for dynamic IO proper.
% =====================================================================
% =====================================================================
% Section 4 --- One kernel, five closures, three financings
% ---------------------------------------------------------------------
% Job: the common kernel (E in passing: one production core deliberately
%      held fixed); row and cell definitions; the "why only three economies"
%      box (Propositions 1-3 in one paragraph, forward reference to Sec. 6
%      and App. A); the presentation rule as one paragraph.
% Mindmap: N02A, N06, N11, N12.
% Status: STUB --- narrative outline v3 (2026-09-20).
% =====================================================================
\section{One kernel, five closures, three financings}\label{sec:kernel}

% TODO: verbal model overview BEFORE the equations (R1.2); state the static
% horizon explicitly; one production core (theory of production E) held fixed
% by design.

We examine how different closures have different effects.
For this we choose five different labor closures and three financing closures that reflect different assumptions about the economy.
The labour closures specify how wages and employment adjust.
\textbf{BF} fixes employment in each sector at its baseline level and lets
sector-specific wages adjust. \textbf{ALPHA} allows labour to move between
sectors at one flexible wage while keeping total employment fixed.
\textbf{BETA} lets sectoral employment respond to each sector's real wage,
deflated by the consumption price index, with a common supply elasticity
but separate wages. \textbf{GAMMA} fixes the common real wage and lets
employment adjust without a cap. \textbf{DELTA} applies GAMMA's wage rule
in the Leontief limit, with fixed proportions in production and consumption;
it is a corner rather than a separate labour-market mechanism.

The financing closures specify how programme-related demand is accommodated
and who bears its cost. \textbf{F1} redirects household spending towards
programme sectors by changing preferences: larger spending shares there are
paid for by smaller shares elsewhere within the household budget, without
additional public purchases. \textbf{F2} adds public investment financed by
a household tax equal to the programme's current-price cost. \textbf{F3}
adds the same public-investment bundle with external financing and no
additional domestic programme tax. This external financing is an accounting
counterpart, not a monetary mechanism.

\subsection{The five labour closures}\label{subsec:labour-closures}
% Formulations: registry/closures.toml; ADR-0020 and ADR-0022.
% Appendix A resets the equation counter; keep these PDF anchors distinct.

\begingroup
\renewcommand*{\theHequation}{labour.\arabic{equation}}
Let $L_i$ denote employment in sector $i$, $\bar{L}_i$ its baseline value,
and $\Lbar = \sum_{i=1}^N \bar{L}_i$ total baseline employment. Sectoral
wages are $w_i$; with mobile labour they equal one common wage $w$.
The real wage is deflated by the consumption price index
\begin{equation}
  \Pi(p) = \left(\sum_{i=1}^N \omega_i p_i^{1-\sigma}\right)^{1/(1-\sigma)},
  \qquad \sum_{i=1}^N \omega_i = 1,
  \label{eq:labour-cpi}
\end{equation}
with calibrated weights $\omega_i$ and the Cobb--Douglas limit at
$\sigma = 1$. Subscript $0$ denotes the baseline. All closures use the
same cost-minimizing labour demand,
\begin{equation}
L_i^{\mathrm{cm}}(p,y_i,w_i)
 = \alpha_i A_i^{\epsilon-1} y_i
   \left(\frac{p_i}{w_i}\right)^{\epsilon}.
\label{eq:closure-labour-demand}
\end{equation}
Here $y_i$ is gross output, $A_i$ productivity and $\alpha_i$ the calibrated
labour share. The elasticity $\epsilon$ governs substitution between labour
and the intermediate composite. The closures differ in the wage and
labour-supply restrictions imposed on this demand.

\paragraph{BF: immobile labour.}
The BF row fixes employment in every sector and solves a sector-specific
wage to clear each labour market:
\begin{equation}
L_i = L_i^{\mathrm{cm}}(p,y_i,w_i) = \bar{L}_i,
\qquad i=1,\ldots,N, \qquad \eta = 0.
\label{eq:closure-bf}
\end{equation}
Thus $\sum_i L_i = \Lbar$, but wages may differ across sectors.
The allocation parameter $\eta$ selects only the endpoints $\{0,1\}$;
its mobile endpoint $\eta = 1$ is ALPHA below, not an additional closure.

\paragraph{ALPHA: mobile labour at full employment.}
Labour reallocates across sectors at one flexible wage, while total
employment remains fixed:
\begin{equation}
L_i = L_i^{\mathrm{cm}}(p,y_i,w),
\qquad \sum_{i=1}^N L_i = \Lbar,
\qquad \eta = 1.
\label{eq:closure-alpha}
\end{equation}

\paragraph{BETA: elastic sectoral labour supply.}
We use the uniform-elasticity sectoral form: each sector has its own wage
and supply curve, with one common labour-supply elasticity $\eta_s$:
\begin{equation}
L_i = L_i^{\mathrm{cm}}(p,y_i,w_i)
 = \bar{L}_i \left[\frac{w_i/\Pi(p)}{w_{i0}/\Pi_0}\right]^{\eta_s},
\qquad \eta_s \geq 0, \qquad i=1,\ldots,N.
\label{eq:closure-beta}
\end{equation}
Equal supply elasticities do not impose equal wages: $N$ separate labour
markets determine $w_i$. Here $\eta_s$ is a labour-supply elasticity, not
the allocation parameter $\eta$; setting $\eta_s = 0$ gives exactly the
immobile BF row. The calibrated real-wage anchors are $w_{i0}/\Pi_0 = 1$.
The deflator remains the common consumption price index, not the sector's
own goods price $p_i$. The supply curves are reduced forms without a
labour--leisure income effect. The headline panels use $\eta_s = 0.5$;
the elasticity ladder and heterogeneous sectoral variants are reported in
Appendix~\ref{app:variants}. The single-wage BETA is retained there only
as a distinct benchmark; its zero-elasticity limit is ALPHA, not BF.

\paragraph{GAMMA: fixed real wage.}
The real wage is pegged at its baseline value, while employment is
endogenous and uncapped:
\begin{equation}
\frac{w}{\Pi(p)} = \bar{w} = 1,
\qquad L_i = L_i^{\mathrm{cm}}(p,y_i,w),
\qquad \eta = 1.
\label{eq:closure-gamma}
\end{equation}
Unlike ALPHA, this closure does not impose $\sum_i L_i = \Lbar$.
It is a two-sided wage peg, not a one-sided wage floor with unemployment
rationing. On the demand-only baseline branch, $\Pi(p)=1$ and the
implemented wage pin is $w=1$.

\paragraph{DELTA: the Leontief corner.}
DELTA retains GAMMA's wage and employment rule while taking the
substitution elasticities to their Leontief limits:
\begin{equation}
\frac{w}{\Pi(p)} = \bar{w} = 1,
\qquad \eta = 1,
\qquad (\theta,\epsilon,\sigma) \longrightarrow (0^+,0^+,0^+).
\label{eq:closure-delta}
\end{equation}
Here $\theta$ governs substitution between intermediate inputs and
$\sigma$ substitution in consumption. DELTA is therefore a corner of
GAMMA with fixed-coefficient technology, rather than an independent
labour-market mechanism.
\endgroup

\subsection{The three financing closures}\label{subsec:financing-closures}
% Formulations: registry/closures.toml; ADR-0002, ADR-0019 and ADR-0020.
% Programme-scale F1 tilt: experiments/run.jl, f1_tilt_weights.
\begingroup
\renewcommand*{\theHequation}{financing.\arabic{equation}}
Each labour closure is crossed with three ways of accommodating the same
programme scale and sectoral incidence. Let $g_i = kG_0\psi_i$ denote the
reference programme bundle, where $G_0$ is its baseline-price value,
$k \geq 0$ scales the programme, and $\psi_i \geq 0$ with
$\sum_i\psi_i = 1$ describes its sectoral incidence. F2 and F3 add this
bundle to final demand; F1 instead uses its scale and incidence to redirect
household spending.

Baseline government purchases $g_i^G$ remain tax-financed in all three
cases. Define the wage bill and their current-price tax cost as
\begin{equation}
W = \sum_{i=1}^N w_i L_i,
\qquad T_G(p) = \sum_{i=1}^N p_i g_i^G.
\label{eq:financing-baseline-budget}
\end{equation}
With a common wage, $W = w\sum_i L_i$. Let $E$ denote after-tax household
resources, including the net external transfer $F$, and let $s$ be the
calibrated saving rate. Household consumption expenditure is
$E^h = (1-s)E$. Gross household purchases $c_i^h$ exhaust this budget:
\begin{equation}
\sum_{i=1}^N p_i c_i^h = E^h,
\qquad c_i^{\mathrm{dom}} = (1-m_i)c_i^h,
\label{eq:financing-household-budget}
\end{equation}
where $m_i$ is the sectoral import margin. Only the domestic content enters
domestic goods-market clearing; likewise, the programme contributes
$(1-m_i)g_i$ under F2 and F3.

\paragraph{F1: budget-neutral preference reallocation.}
F1 changes the household's CES weights without adding a public-demand
bundle. For positive preference shifters $d_i$, the weights and gross
household demand become
\begin{align}
\widetilde{\omega}_i
 &= \frac{\omega_i d_i}{\sum_j\omega_j d_j},
\qquad \sum_i\widetilde{\omega}_i = 1,
\label{eq:closure-f1-weights}\\
c_i^h
 &= E^h\frac{\widetilde{\omega}_i p_i^{-\sigma}}
 {\sum_j\widetilde{\omega}_j p_j^{1-\sigma}},
\qquad E = W-T_G(p)+F.
\label{eq:closure-f1}
\end{align}
The programme-scale tilt is $d_i = 1+kG_0\psi_i^+/c_{i0}^h$ in sectors
with positive baseline household purchases, and $d_i=1$ otherwise;
$\psi^+$ is the programme incidence restricted to those sectors and
renormalized to sum to one. The CES normalizer preserves
\eqref{eq:financing-household-budget} at every price vector. Thus increased
spending shares in programme sectors are paid for by reduced shares
elsewhere, conditional on the household budget; equilibrium income need
not remain unchanged. F1 is a demand-composition experiment, not an
autonomous investment injection or multiplier.

\paragraph{F2: tax-financed public investment.}
F2 adds the real programme bundle $g$ and levies an additional household
tax equal to its value at equilibrium prices:
\begin{equation}
T(p) = \sum_{i=1}^N p_i g_i,
\qquad E = W-T_G(p)-T(p)+F,
\qquad B_{\mathrm{gov}}=0.
\label{eq:closure-f2}
\end{equation}
Total domestic tax revenue $T_G(p)+T(p)$ therefore balances the purchases
$g^G+g$, priced at current rather than baseline prices. The tax is
implemented as the endogenous wage-income rate
$\tau=[T_G(p)+T(p)]/W$; with one representative household its burden is
effectively lump-sum. The programme has an explicit domestic budget
counterparty and reduces household resources directly by $T(p)$.

\paragraph{F3: externally financed public investment.}
F3 adds the same real bundle but imposes no additional domestic programme
tax. Its current-price cost is booked as external programme financing:
\begin{equation}
B_{\mathrm{gov}} = \sum_{i=1}^N p_i g_i,
\qquad E = W-T_G(p)+F.
\label{eq:closure-f3}
\end{equation}
The booking $B_{\mathrm{gov}}$ is distinct from the net external transfer
$F$ entering household resources after tax. The latter is solved
endogenously in BF, ALPHA and both BETA forms, but is fixed at zero in
GAMMA and DELTA. With no unfinanced demand injection, all closures clear
all $N$ goods markets and satisfy the external-account identity
\begin{equation}
S + T_{\mathrm{int}} + M - (I+X) = F+B_{\mathrm{gov}},
\qquad S=sE,
\label{eq:financing-external-account}
\end{equation}
where $I$ is non-programme investment, $X$ exports, $M$ total imports
(including intermediate and programme imports), and $T_{\mathrm{int}}$
the separately booked product-tax leak on intermediate use; all are
current-price values. Here $B_{\mathrm{gov}}=0$ under F1 and F2, so the net
external position is $F+B_{\mathrm{gov}}$, not the programme cost alone.
F3 is an accounting open-economy closure, not a monetary mechanism: the
static model cannot distinguish external borrowing from a transfer and
does not model interest payments, repayment or exchange-rate adjustment.
\endgroup

\subsection{Why only three economies}\label{subsec:why-three}
% TODO: the equivalence box --- Propositions 1-3 in one paragraph
% (demand-only scope, the A = 1 branch the numeraire selects); forward
% reference to Sec. 6 and App. A.

\subsection{Presentation rule}\label{subsec:presentation-rule}
% TODO: ONE paragraph. Columns F1, F2, F3; shade F2 = F3 in the BF and
% ALPHA/BETA rows only; state the neutrality condition (F endogenous); flag
% the fixed-wage row as the financing-sensitive case; restate the collapse as
% fifteen cells -> seven economies (NOT three; three is a count of labour
% classes). Reference paper/framing_gamma.tex. Keep this as one paragraph ---
% do not leak it as a debate.

% ===== CARRIED OVER FROM AN OLD DRAFT =============================
% [devised for an old draft; not in the validated submission]
% source: revised_manuscript/chapters/ch03_cge_model.tex (commit 41144dd)
% =================================================================
% Eventually, the question emerges how to link such a setup with the Input-Output tables introduced in the beginning of section \ref{sec:IO_models}. In practice, such a link necessitates that $\Omega$ plays a role in the exact specification of the production function and/or consumption function. For a Cobb-Douglas function, as the most widely used specification, the cost shares contained in $\Omega$ are interpreted as output elasticities and, correspondingly, used to specify the relative contribution of factors and intermediates to net output of a given sector, which translates into
% and choose the CES consumer price index as the numeraire:$$
% \text{CPI} = \left( \sum_{u=1}^n \omega_{0u} p_u^{\theta  -1} \right)^{\frac{1}{\theta - 1}}.
% $$
% Calculating GDP directly using these post-shock prices would yield a nominal GDP measure. To obtain real GDP, one typically employs a price index such as the Laspeyres or Paasche index. We use the Laspeyres index, so that GDP is computed as
% $$
% \text{GDP} = \frac{\sum_{s=1}^n p_s \post c_s}{\sum_{s=1}^n p_s  c_s}
% $$
% which measures the change in consumption valued at pre-shock prices, providing a meaningful comparison of real economic activity before and after the shock.
% While this basic formal setup obviously differs from the Leontief model presented in Section~\ref{sec:leontief} in many details,
% a major practical difference in applying these models also emerges from the underlying paradigmatic perspectives:
% loosely spoken, the Leontief model follows a 'demand-side view', where production structures are assumed constant and final demands change exogenously. In contrast, general equilibrium approaches will start from a 'supply-side view', where economic changes arise from changes in the exogenous parameters -- like production functions, technology or endowments -- and demand will respond.
% NOTE: the closing 'demand-side vs supply-side view' sentence was rejected by Reviewer 2 (R2.7) as nonsensical in GE; keep only if reworked.

% =====================================================================
% Section 5 --- Application: the programme through the matrix
% ---------------------------------------------------------------------
% Job: multipliers, price responses, external position across cells; the
%      robust-vs-closure-sensitive table; decision heuristics for
%      transformation modellers. THIS IS THE HEADLINE THE INTRO PROMISES.
% Mindmap: N06, N05, N10.
% Status: STUB --- narrative outline v3 (2026-09-20).
% =====================================================================
\section{Application: the programme through the matrix}\label{sec:application}

% TODO: run the green-investment programme through the 5x3 matrix.
% Every table cites run_ids from matrix_5x3_v10 (matrix + 18 sectoral cells)
% and supply_etas_* (189 cells; s1 the identification vehicle).
% TABLES (linked, revision/tables/):
%   - matrix_5x3_v10_flows.md  -> the 15 matrix cells + 18 sectoral cells (Tables 1-4)
%   - supply_etas_flows.md     -> the supply arm s1 (identification vehicle)
%   - supply_etas_prog_flows.md / supply_etas_unif_flows.md -> the other two designs
% Port each to LaTeX at its point of use.

\subsection{Aggregate results}\label{subsec:aggregate}
% TODO: multipliers and price responses; the wage-regime vs elasticity
% channels (R1.7). Real GDP is the income-side index, flat by construction
% under demand-only shocks (ADR-0018); consumption is welfare. Do NOT repeat
% the old "virtually zero to ~1.5 %" GDP sentence (retired framing).

\subsection{Sectoral results}\label{subsec:sectoral}
% TODO: sectoral detail; the ladder and rigid-group cells.

\subsection{The financing wedge}\label{subsec:financing-wedge}
% TODO: bracket results (F1/F2/F3); the fixed-wage row is where financing
% bites (GAMMA/DELTA F2 vs F3: -1.53 % vs +2.26 % consumption).

\subsection{Robust vs closure-sensitive}\label{subsec:robust}
% TODO: the decision table for transformation modellers; which conclusions
% survive every closure, which flip.

% ===== CARRIED OVER FROM AN OLD DRAFT =============================
% [devised for an old draft; not in the validated submission]
% source: revised_manuscript/chapters/ch05_ces_functions.tex (commit 41144dd)
% =================================================================
% An overall inspection of the results show that in this comparison the standard Leontief approach delivers the largest GDP-estimates -- on this basis GDP is expected to grow by about 1.5\%, which is equivalent to a multiplier of 1.15 given that the initial investment impulse was specified to be 40.337 bn. €, about 1.3\% of GDP. In contrast, the Cobb-Douglas specification predicts a small reduction in GDP that results from efficiency-losses arising from the increase in green investment. These losses are due to the fact that the imposed green investment leads to a forced reallocation of intermediates and production factors towards the expanding sectors, which alters an already optimal allocation and, hence, induces inefficiencies. In turn, real output decreases slightly and the major macroeconomic impact of a green investment demand shock in this framework is to induce sector-specific inflationary pressures.
% Finally, by assessing changes in sectoral production after the demand shock in Figure \ref{fig:sectoral-changes}, we find that the output adjustment approach will lead to an expansion of almost all sectors, while the general equilibrium approach leads to a reduction of gross output in the large majority of sectors, that are not directly affected by the demand shock. Moreover, the shock also depresses gross output and final consumption associated with two of the shocked sectors (in line with Figure \ref{fig:comparison-sec4} above, which focuses on final demand only). This reproduces the finding from section \ref{subsec:CGE_elasticity} that a demand shock will depress real GDP on a sectoral level as the (comparatively lower) average expansion of shocked sectors does under no conditions fully compensate for the reduction in the production of unaffected sectors. Notwithstanding this overall pattern, the complex substitution dynamics implied the CGE approach lead to exceptional outcomes in specific sectors. The most obvious idiosyncratic outcome occurs in two sectors that are not shocked –– \textit{building construction works (sector 33 )} and \textit{civil engineering works (sector 34)} --, but, rather, massively expand their intermediate goods production at the expense of final outputs, which represents an over-proportional intersectoral crowding-out effect. In other words, intermediary goods from these sectors are in extremely high demand due to the enforced shift in final expenditures. For both of these sectors the expansion in intermediate goods production is so strong, that these sectors expand more strongly than in the Leontief case and thereby partially compensate for the lack of freely available workers in shocked sectors.


% =====================================================================
% Section 6 --- Beyond demand-only shocks
% ---------------------------------------------------------------------
% Job: what simplifies and where variety lives --- the equivalences and
%      their scope (demand-only, the A = 1 branch the numeraire selects);
%      the sectoral family and its exact nesting (the rigid corner IS the
%      eta = 0 endpoint); the executed ladder and rigid-group cells; the
%      supply-shock arm as identification vehicle; the no-go corollary; BF
%      repositioned here as the frictionless pole.
% Mindmap: N03, N13, N02D.
% Status: STUB --- narrative outline v3 (2026-09-20).
% =====================================================================
\section{Beyond demand-only shocks}\label{sec:beyond-demand-only}

% TODO: scope the equivalences airtight; give a clean supply-shock
% counterexample and use it to motivate this section rather than hide the
% boundary. Reference paper/framing_gamma.tex.

\subsection{The equivalences and their scope}\label{subsec:equivalences}
% TODO: Lemma 1 (price-block demand-invariance), Lemma 2 (nonsubstitution,
% Samuelson 1951), Propositions 1-3; financing neutrality; full proofs in
% App. A.

\subsection{The sectoral family}\label{subsec:sectoral-family}
% TODO: L^cm_i = Lbar_i (w_i/Pi)^{eta_s,i}; the rigid corner is exactly the
% eta = 0 endpoint (Prop. 3); the executed ladder (BETA-F2-etas{025,05,1,2} =
% 0.152897 / 0.105666 / 0.065422 / 0.037171) and the rigid-group cells.

\subsection{The supply-shock arm}\label{subsec:supply-arm}
% TABLE (linked, revision/tables/supply_etas_flows.md): the s1 identification
% vehicle; supply_etas_prog_flows.md / supply_etas_unif_flows.md for the other
% two designs (189 cells total).
% TODO: the identification vehicle (supply_etas_s1); 189 cells across three
% designs; separates the ladder rungs that demand-only shocks cannot.

\subsection{The no-go corollary and the BF pole}\label{subsec:no-go}
% TODO: nominal wage rules are vacuous in both gauges --- the rigidity that
% matters must live in sectoral labour markets. BF repositioned as the
% frictionless pole (peer-modesty: late introduction, no vantage status).
% ===== CARRIED OVER FROM AN OLD DRAFT =============================
% [devised for an old draft; not in the validated submission]
% source: revised_manuscript/chapters/ch05_ces_functions.tex (commit 41144dd)
% =================================================================
% With regard to how changes in assumptions on substitution elasticities impact on final outcomes Figure \ref{fig:elasticity-gradient} indicates that real GDP is typically increasing in all three elasticities employed -- the higher the chosen elasticities, the higher are the respective GDP-estimates. This finding reproduces the intuition that easier substitution will lead to higher (aggregate) output and, hence, CES-based estimates are consistently lower than their Cobb-Douglas counterparts.
% 
% Against this backdrop we concede that the introduction of Leontief-like properties in the production function helps bridging paradigmatic cleavages on a level of conceptual foundations -- e.g. with respect to the question how to adequately represent production processes. However, we find that introducing such a conceptual shift in neoclassical IO-models will actually serve to increase the gap between both approaches in terms of predicted outcomes. This result is not too surprising when taking into account that the assumption of a Leontief technology implies less flexibility as compared to a Cobb-Douglas technology. Hence, in the context of assumptions on substitutability neoclassical approaches are more optimistic than the heterodox archetype: viewed in isolation greater substitutability comes with higher prospects for growth in case of demand shocks and, correspondingly, leads to a more positive assessment of fiscal intervention in terms of forced green investment. However, this relatively more optimistic stance does in no way compensate for the fact that the estimated growth effects remain negative, which is effectively driven by assumptions on capacity constraints. However, before turning to this key aspect in section \ref{sec:labor} further below, we first inspect structural differences between model outcomes on a more granular level.


% ===== CARRIED OVER FROM AN OLD DRAFT =============================
% [devised for an old draft; not in the validated submission]
% source: revised_manuscript/chapters/ch06_labor_slack.tex (commit 41144dd)
% =================================================================
% Nonetheless, on a general level this example again illustrates how the introduction of labor slack can lead to more optimistic results in terms of GDP growth. The key requirement is the willingness to assume that such an expansion of the labor force is indeed a plausible reaction -- either due to labor market considerations or as a general response to increased effective demand. While this feature is not truly surprising it shows how modifying dominant convictions on the role and modeling of capacity constraints indeed opens up a shared realm, in which the simplistic juxtaposition of an 'expansionary' Leontief-approach and a 'contractionary' CGE-perspective becomes subject to a more nuanced perspective that allows for pragmatically interpreting labor slack as a moderating variable to explore some shared ground between these two competing visions of the economy.


% =====================================================================
% Section 7 --- Discussion
% ---------------------------------------------------------------------
% Job: positioning against BF and the production-network literature (pole,
%      not foil); FIDELIO and structuralist macro; limitations; open doors
%      (capacity, markup) acknowledged as future work.
% Mindmap: N03, N02D.
% Status: STUB --- narrative outline v3 (2026-09-20).
% =====================================================================
\section{Discussion}\label{sec:discussion}

% TODO: BF (2019, 2021, 2022) and the production-network literature --- BF is
% the pole, not the foil. FIDELIO (Rocchi et al.) and structuralist macro
% (Taylor 1990). Acemoglu, Akcigit & Kerr (2015) on networks and the macro.
% Limitations: one-factor aggregation, static horizon, open-economy
% treatment, welfare metric (dispersion-blind). Open doors: capacity
% (Kaldor-Verdoorn), markup pricing, skill classes (R2.14), dynamic CGE.
% =====================================================================
% Section 8 --- Conclusion
% ---------------------------------------------------------------------
% Job: return to the thesis --- closure choice is first-order, and the
%      matrix tells you where it is and is not.
% Mindmap: ---
% Status: STUB --- narrative outline v3 (2026-09-20).
% =====================================================================
\section{Conclusion}\label{sec:conclusion}

% TODO: return to the thesis sentence (the transformation short run is doubly
% rigid; BF the frictionless pole). Practical closure guidance for
% transformation modellers. NO bridge-discovery claims (R2.1/R2.3 are
% answered by measurement, not by a new bridge). Keep the Kuhnian frame to
% the minimum (Kapeller 2013; Kuhn 1962).

\section*{Data availability}
The input--output data for Germany employed in this paper is publicly available
(see the Destatis input--output accounts cited in Appendix~\ref{app:provenance}).

\section*{Funding}
This paper has benefited from funding provided by the \textit{Hans B\"ockler
Foundation} under grant number 2021-544-2, by the \textit{Linz Institute for
Transformational Change (LIFT\_C)} under grant number 2022-1-114 as well as by
\textit{Dezernat Zukunft}.

\section*{Conflict of interest}
The authors report no conflicts of interest associated with the content of this paper.

\medskip
\printbibliography
\ifendfloats\processdelayedfloats\fi

\appendix
\begin{refsection}

\setcounter{section}{0}
\setcounter{table}{0}
\setcounter{equation}{0}
\setcounter{figure}{0}
\ifendfloats
  \setcounter{postfigure}{0}
  \setcounter{posttable}{0}
  \renewcommand{\theposttable}{A\arabic{posttable}}
  \renewcommand*{\theHposttable}{appendix.posttable.\arabic{posttable}}
  \renewcommand{\thepostfigure}{A\arabic{postfigure}}
\fi
\renewcommand{\thetable}{A\arabic{table}}
\renewcommand*{\theHtable}{\thetable}
\renewcommand{\thefigure}{A\arabic{figure}}
\renewcommand*{\theHfigure}{\thefigure}
\renewcommand{\theequation}{A\arabic{equation}}

% =====================================================================
% Appendix A --- Full proofs
% ---------------------------------------------------------------------
% Ported from paper/to_evaluate/equivalence.tex (v4, evidence layer
% matrix_5x3_v10, 2026-09-20). Standalone front matter (title, abstract,
% version block, revision log) dropped; sections demoted to subsections;
% revision-colour wrappers stripped; the two labels colliding with the
% manuscript (sec:kernel, sec:conclusion) retargeted to app:eq*.
% =====================================================================
\section{Full proofs}\label{app:proofs}

\noindent This appendix reproduces, in full, the ex ante equivalence results
of the revision (source note: \texttt{equivalence.tex}, v4). It states the
common kernel, derives the demand-only reduction, proves the exact
equivalences of the closure rows and of the sectoral family's rigid corner,
isolates what does differ across the rows, and derives the scope conditions
under which the equivalences break.

\subsection{Purpose and summary}

Reviewer 2's critique turns on the exact conditions under which the closure
rows of the evaluation matrix are or are not equivalent
(\texttt{docs/DOCS\_ASSESSMENT.md}, section 4.2). The executed
\texttt{matrix\_5x3\_v10} generation (ADR-0019 together
with the ADR-0020 option C endpoint; 33/33 cells, run ids
\texttt{matrix\_5x3-v10-*}) measures two coincidences:

\begin{itemize}
\item \textbf{ALPHA $\equiv$ BETA.} The elastic-supply row
      ($\eta_s = 0.5$) reports every metric of the full-employment row
      identically: real GDP, the welfare index, employment, the external
      account --- to all eight printed decimals, with the identity wedge at
      $\sim 10^{-16}$. A direct comparison of stored solutions measured
      $\max|\Delta y| = 6.2\times 10^{-17}$ (session of 2026-09-17).
      The v6 generation reproduces the v5 ALPHA/BETA
      cells to $6.7\times 10^{-16}$, so the coincidence survives the endpoint
      change.
\item \textbf{GAMMA $\equiv$ DELTA.} The Leontief-corner row
      ($( \theta_\delta, \epsilon_\delta, \sigma_\delta ) = (10^{-4},
      10^{-4}, 10^{-4})$) reports every metric of the fixed-wage row at the
      baseline CES elasticities $(\theta,\epsilon,\sigma) = (0.5,0.5,0.9)$
      identically, in all three financing columns; the flow-table wedges
      agree at the residual level ($\sim 10^{-12}$). The analytic
      \texttt{leontief\_multiplier} test reproduces the nonlinear DELTA
      solution with relative error $0.0$ (recorded in
      \texttt{docs/DOCS\_ASSESSMENT.md} section 4.2).
\end{itemize}

Sections~\ref{app:eqkernel}--\ref{sec:reduction} set up the common kernel and
derive the demand-only reduction. Section~\ref{sec:equiv} states and proves
the two equivalences and classifies them.
Section~\ref{sec:sectoral} records a third exact
equivalence, of a different kind: the rigid corner of the sectoral-labour-market
family is the $\eta = 0$ endpoint itself, and the family interpolates between the
matrix's two rigidity rows. Section~\ref{sec:gauge} isolates
what \emph{does} differ across the rows (the scalar closure regimes).
Section~\ref{sec:scope} derives the scope conditions: which model changes
break the equivalences --- the menu behind the demand-sensitive-price program
--- and which cannot.

Throughout, ``demand-only shock'' means the matrix design: a programme bundle
$g = G_0\,\psi$ ($\sum_i \psi_i = 1$; $G_0 = 40\,300$
EUR m, ADR-0024) entering final demand under one of the three financing
closures, with the supply shock at $A \equiv 1$ (unit productivity).

\subsection{The common kernel}\label{app:eqkernel}

Let $N$ be the number of sectors. Prices $p_i$, wage $w$, gross output
$y_i$, employment $L_i$. Data-calibrated objects: factor shares $\alpha_i$
(labour), domestic bill coefficients $a_u = A^{bill}_u/\lambda_u$ with
input-share matrix $\Omega$ ($\Omega_{iu}$ = the share of input
$u$ in sector $i$'s input bundle; rows sum to one), CPI weights $\omega_i$ ($\sum_i \omega_i = 1$),
sector import margins $m_i$, saving rate $s$, exogenous government demand
$g^G$, investment $inv$, exports $X$ (no margin), the programme bundle $g$,
and the labour bar $\Lbar = \sum_i \bar{L}_i = 1$ (baseline income units,
GDP $=1$). Elasticities: intermediate substitution $\theta$,
labour--intermediate substitution $\epsilon$, consumption substitution
$\sigma$, allocation $\eta \in \{0,1\}$ (ADR-0010), supply elasticity
$\eta_s \geq 0$ (BETA only).

\paragraph{Price block.} The CES intermediate price index and unit cost are
\begin{align}
P_i(p) &= \Big( \textstyle\sum_u \Omega_{iu}\, p_u^{\,1-\theta}
   \Big)^{\frac{1}{1-\theta}}, \label{eq:ip}\\
c_i(p,w;A) &= \Big( A_i^{\,\epsilon-1}\big( \alpha_i\, w^{1-\epsilon}
   + (1-\alpha_i)\, P_i^{\,1-\epsilon} \big) \Big)^{\frac{1}{1-\epsilon}},
   \label{eq:cost}
\end{align}
with the Cobb--Douglas limit $c_i = w^{\alpha_i} P_i^{\,1-\alpha_i}/A_i$ at
$\epsilon = 1$. Zero profit requires
\begin{equation}
p_i = c_i(p,w;A) \qquad \forall\, i. \label{eq:zp}
\end{equation}
At the calibration point $(p,w,A) = (\one,1,\one)$ the unit cost equals one
exactly, because $\sum_u \Omega_{iu} = 1$ makes $P_i = 1$; this is the
calibration identity used throughout.

\paragraph{Quantity block.} With full cost-minimizing allocation
($\eta = 1$) the marginal-product condition inverts to
\begin{equation}
L_i = \big( p_i\, A_i^{(\epsilon-1)/\epsilon}\, \alpha_i^{1/\epsilon}\,
      y_i^{1/\epsilon} / w \big)^{\epsilon}
    = p_i^{\epsilon}\, A_i^{\epsilon-1}\, \alpha_i\, y_i / w^{\epsilon},
\label{eq:labdemand}
\end{equation}
At $\eta = 0$ the allocation is the frozen baseline, $L_i = \bar{L}_i$;
under ADR-0020 option C the wage $w$ in
\eqref{eq:labdemand} is then a per-sector $w_i$, solved from
$L^{cm}_i(p, y_i, w_i) = \bar{L}_i$ jointly with \eqref{eq:zp}.
Intermediate demand is charged with the \emph{domestic}
bill (ADR-0012):
\begin{equation}
int_i = p_i^{-\theta} \sum_u \Omega_{ui}\,
        p_u^{\epsilon} A_u^{\epsilon-1} P_u^{\theta-\epsilon}\, a_u\, y_u.
\label{eq:int}
\end{equation}
Household demand is CES with CPI deflation and the financing weights
$ds$ ($ds = 1$ except under F1, where $ds = d$, the preference tilt):
\begin{equation}
c_i = (1-s)(1-m_i)\,\omega_i\, ds_i\; E\; p_i^{-\sigma} \Big/
      \textstyle\sum_j \omega_j ds_j\, p_j^{1-\sigma},
\qquad
E = (1-\tau)\, w \textstyle\sum_i L_i + F,
\label{eq:cons}
\end{equation}
with the lump-sum tax rate $\tau = p\!\cdot\! g^G / (w \sum_i L_i)$ under
F1/F3 and $\tau = p\!\cdot\!(g^G + g)/(w \sum_i L_i)$ under F2, and $F$ the
net external transfer entering expenditure after tax (ADR-0019).
At $\eta = 0$ the wage bill is
$\sum_i w_i \bar{L}_i$ in \eqref{eq:cons}. Saving
$sE$ leaks; the import content is charged margin by margin. Goods-market
clearing reads
\begin{equation}
y_i = int_i + c_i + (1-m_i)\big( g^G_i + inv_i + g_i \big) + X_i,
\label{eq:clearing}
\end{equation}
where $g_i$ appears only under the additive financings F2/F3
($g_i = 0$ under F1; legacy manna is zero in the matrix).

\paragraph{The four regimes.} The mobile systems solve for
$X = (p, y, w, F) \in \mathbb{R}^{2N+2}$ with the numeraire
$\Pi(p) = 1$, where
\begin{equation}
\Pi(p) = \Big( \textstyle\sum_i \omega_i\, p_i^{1-\sigma}
   \Big)^{\frac{1}{1-\sigma}},
\label{eq:cpi}
\end{equation}
and close with the labour equation
\begin{align}
\text{ALPHA:}\quad & \textstyle\sum_i L_i(p,y,w) = \Lbar,
\label{eq:alpha}\\
\text{BETA:}\quad & \textstyle\sum_i L_i(p,y,w)
   = \Lbar \Big[ \frac{w/\Pi(p)}{w_0/\Pi_0} \Big]^{\eta_s},
\qquad w_0/\Pi_0 = 1.
\label{eq:beta}
\end{align}
The fixed-wage systems solve for $X = (p, y) \in \mathbb{R}^{2N}$ with the
wage pinned $w \equiv 1$ and $F \equiv 0$; employment is a post-solve
outcome. GAMMA runs them at the baseline elasticities, DELTA at the Leontief
corner $(\theta_\delta,\epsilon_\delta,\sigma_\delta) = (10^{-4},10^{-4},
10^{-4})$. BF is the sectoral-wage system at
$\eta = 0$ (ADR-0020 option C, superseding the ADR-0019 $F = 0$ pin): the
allocation stays frozen at the baseline, $L_i = \bar{L}_i$, and $N$ sectoral
wages $w_i$ are solved from each sector's own first-order condition
$L^{cm}_i(p, y_i, w_i) = \bar{L}_i$ jointly with zero profit, with unknowns
$[\,p;\; y;\; w(1{:}N);\; F\,]$ and the external transfer $F$ free. All rows
in the matrix use $\eta = 1$ except BF.

\subsection{The demand-only reduction}\label{sec:reduction}

\begin{lemma}[Price block demand-invariance]\label{lem:price}
The zero-profit system \eqref{eq:zp} is homogeneous of degree one in
$(p,w)$ jointly and contains no demand variable. Consequently, for any
technology $A$, every solution satisfies $p = w\,\pi(A)$, where
$\pi(A) \in \mathbb{R}^N_{++}$ solves $c_i(\pi,1;A) = \pi_i$ for all $i$;
relative prices and the ratio $p/w$ are functions of $A$ alone. Under the
CPI numeraire, $w = 1/\Pi(\pi(A))$ and $p = \pi(A)/\Pi(\pi(A))$; under the
fixed wage, $w = 1$ and $p = \pi(A)$. At $A \equiv \one$ the calibration
identity gives $\pi = \one$, hence in both gauges
\begin{equation}
p = \one, \qquad w = 1, \qquad \Pi = 1,
\end{equation}
independently of the demand vector, of the financing closure, and of the
shock size. The lemma is a statement about the
\emph{single-wage} price block: with a per-sector wage vector (the
$\eta = 0$ endpoint, ADR-0020) homogeneity in $(p, w)$ no longer holds
sector by sector and the argument does not apply.
The statement is scoped to the branch the numeraire
selects: $\pi(A)$ denotes the normalised solution the kernel's solves return,
and the equivalences below compare the systems on that branch (the CES price
system at $\theta<1$ with self-loops admits spiral branches in principle; no
executed cell visits one).
\end{lemma}

\begin{proof}
Homogeneity: $c_i(tp, tw; A) = t\, c_i(p,w;A)$ for $t>0$, by inspection of
\eqref{eq:cost}--\eqref{eq:ip}. Write $q = p/w$. Then \eqref{eq:zp} is
equivalent to $q_i = c_i(q,1;A)$ for all $i$, i.e.\ $q = \pi(A)$, and
the solution branch of \eqref{eq:zp} selected by the
numeraire is the ray $\{(w\,\pi(A), w)\}$. The
CPI \eqref{eq:cpi} is homogeneous of degree one in $p$, so
$\Pi(w\,\pi) = w\,\Pi(\pi)$; the numeraire $\Pi = 1$ or the pin $w = 1$
selects the scale. At $A \equiv \one$, $\pi = \one$ solves
\eqref{eq:zp} at $w = 1$ because $c_i(\one,1;\one) = 1$ (the calibration
identity), and $\Pi(\one) = 1$ because $\sum_i \omega_i = 1$.
\end{proof}

\begin{corollary}[Real-wage invariance]\label{cor:realwage}
At every solution of the shared block,
\begin{equation}
\frac{w}{\Pi(p)} = \frac{1}{\Pi(\pi(A))},
\end{equation}
a function of technology alone. At $A \equiv \one$ this equals $1 =
w_0/\Pi_0$: the equilibrium real wage sits at its anchor for every demand
vector, every financing closure and every shock magnitude. (Measured in the
matrix: the GDP deflator reads $1.000000$ in the
twelve $\eta = 1$ cells, and a demand
shock of ten times the programme size moves $\max|p-\one|$ by
$2.7\times10^{-15}$; \texttt{docs/ETAs.md}.)
\end{corollary}

\begin{lemma}[Elasticity-free affine quantity block]\label{lem:affine}
Fix $\eta = 1$ and $A \equiv \one$, and evaluate the quantity block at the
demand-only prices $(p,w) = (\one,1)$. Then for \emph{every} value of
$(\theta,\epsilon,\sigma)$:
\begin{align}
int &= \Omega^{\!\top}\mathrm{diag}(a)\, y, &
L &= \mathrm{diag}(\alpha)\, y, &
\sum_i L_i &= \alpha^{\!\top} y,
\label{eq:leontiefblocks}\\
c &= (1-s)\,\mathrm{diag}(1-m)\,\tilde\omega(ds)\; E, &
E &= \alpha^{\!\top} y - \textstyle\sum g^G + F
   \;\; (\text{F1/F3}), \nonumber
\end{align}
where $\tilde\omega(ds) = (\omega \odot ds)\,/\, \one^{\!\top}(\omega \odot
ds)$ and $E = \alpha^{\!\top} y - \sum g^G - \sum g$ under F2. Hence the
clearing system \eqref{eq:clearing} is exactly affine,
\begin{equation}
y = G\, y + b(F),
\qquad
G = \Omega^{\!\top}\mathrm{diag}(a)
  + (1-s)\,\mathrm{diag}(1-m)\,\tilde\omega(ds)\,\alpha^{\!\top},
\label{eq:affine}
\end{equation}
\begin{equation}
b(F) = (1-m)\odot\big( g^G + inv + g \big) + X
     + (1-s)\,\mathrm{diag}(1-m)\,\tilde\omega(ds)
       \big( F - \textstyle\sum g^G - \mathbf{1}_{F2}\textstyle\sum g \big),
\label{eq:bvec}
\end{equation}
with $g$ present only under the additive financings
F2/F3 and $\mathbf{1}_{F2}$ selecting the tax-financed column (the bracket is
$F-\sum g^G$ under F1/F3 and $F-\sum g^G-\sum g$ under F2).
No elasticity parameter enters $G$ or $b$. At
$\eta = 0$ the aggregate is still pinned, $\sum_i L_i = \Lbar$, but the
unit-price reduction below describes only the \emph{retired} common-wage
pin: there $E$ did not depend on $y$ and the household round dropped out,
$G_{BF} = \Omega^{\!\top}\mathrm{diag}(a)$. Under the ADR-0020 sectoral-wage
closure the wages $w_i$ move with the demand composition, the prices leave
the unit point, and this affine form does not describe the $\eta = 0$ row.
\end{lemma}

\begin{proof}
At $p = w = A = \one$: $P_u = 1$ in \eqref{eq:int}, so all elasticity
factors $p^{\epsilon}$, $P^{\theta-\epsilon}$, $p^{-\theta}$ equal $1$ and
\eqref{eq:int} reduces to the fixed-coefficient round
$\sum_u \Omega_{ui} a_u y_u$; \eqref{eq:labdemand} reduces to
$L_i = \alpha_i y_i$ since every exponent collapses at unit prices; in
\eqref{eq:cons} the terms $p_i^{-\sigma}$ and the normalizer reduce to
$1$ and $\sum_j \omega_j ds_j$, giving the fixed spending weights
$\tilde\omega(ds)$; and the tax rate $\tau = \sum g^G / \alpha^{\!\top}y$
turns $E$ into the affine form shown. Collecting terms in
\eqref{eq:clearing} gives \eqref{eq:affine}--\eqref{eq:bvec}: the
intermediate round contributes the columns
$\sum_i \Omega_{ui} a_u = a_u$ (row sums of $\Omega$ equal $1$) and the
household round contributes $(1-s)\sum_i (1-m_i)\tilde\omega_i(ds)\,
\alpha_u$ per unit of $y_u$.
\end{proof}

\begin{remark}[Nonsubstitution]\label{rem:samuelson}
Lemma~\ref{lem:price} and Lemma~\ref{lem:affine} are the production side of
the Samuelson nonsubstitution theorem: with one primary factor, constant
returns and no joint production, the cost-minimizing technique --- and hence
relative prices --- is independent of the demand vector, and final demand
allocates quantities only (Samuelson 1951). The kernel's nested-CES
structure is a smooth parametrisation of that corner; the demand-only matrix
therefore inherits the theorem's rigidity, and this is what
\texttt{docs/DOCS\_ASSESSMENT.md} section 4.2 records as the ``model-class
property'' behind the coincidences. The IO endpoint makes the rigidity
literal --- the Robinson (2006) fixed-price multiplier --- but, as
Proposition~\ref{prop:gd} shows, it is not an approximation:
the same candidate equilibrium $(p,y)=(\one,y^{\ast})$
solves the fixed-wage system at every elasticity triple in the fixed-wage
gauge.
\end{remark}

\subsection{The two equivalences}\label{sec:equiv}

\begin{proposition}[ALPHA $\equiv$ BETA]\label{prop:ab}
Let $A \equiv \one$ and $\eta = 1$. The ALPHA system
\eqref{eq:zp},\eqref{eq:clearing},\eqref{eq:alpha},$\Pi = 1$ and the BETA
system \eqref{eq:zp},\eqref{eq:clearing},\eqref{eq:beta},$\Pi = 1$ have
identical solution sets, for every $\eta_s \geq 0$ and every financing
closure.
\end{proposition}

\begin{proof}
The two systems differ in exactly one equation. By
Lemma~\ref{lem:price} and Corollary~\ref{cor:realwage}, every solution of
the shared block satisfies $w/\Pi = 1 = w_0/\Pi_0$, so BETA's labour
equation \eqref{eq:beta} reduces to $\sum_i L_i = \Lbar$, which is
\eqref{eq:alpha}. Conversely, any solution of the ALPHA system satisfies
BETA's equation by the same substitution. The systems are therefore
equivalent, equation by equation.
\end{proof}

\begin{proposition}[GAMMA $\equiv$ DELTA]\label{prop:gd}
Let $A \equiv \one$, $\eta = 1$, $F \equiv 0$, $w \equiv 1$. Then
$(p,y) = (\one,\, y^{\ast})$ with
$y^{\ast} = (I - G)^{-1} b(0)$ solves the fixed-wage system at
\emph{every} elasticity triple $(\theta,\epsilon,\sigma)$; the employment
outcome is $L^{\ast} = \alpha^{\!\top} y^{\ast}$, the consumption block is
$c^{\ast} = (1-s)\,\mathrm{diag}(1-m)\,\tilde\omega(ds)\,
(L^{\ast} - \sum g^G - \sum g\,\mathbf{1}_{F2})$, where
$\mathbf{1}_{F2}$ selects the tax-financed column, and the block
decomposition \eqref{eq:leontiefblocks} holds verbatim. In particular GAMMA (baseline
elasticities) and DELTA (the Leontief corner) share this solution, and the
analytic system \eqref{eq:affine}--\eqref{eq:bvec} ---
\texttt{leontief\_multiplier} in \texttt{src/closures/labor/labor.jl} ---
computes it exactly.
\end{proposition}

\begin{proof}
Existence and form: at $p = \one$ the price block holds by the calibration
identity (Lemma~\ref{lem:price}), and by Lemma~\ref{lem:affine} the
clearing block at $p = \one$ is the affine system, of which
$y^{\ast} = (I-G)^{-1}b(0)$ is the solution. Uniqueness on the admitted
branch: the scale-determinacy guard (ADR-0017) verifies at the reference
prices that $(I - G)$ is nonsingular --- contraction if
$\max_u \mathrm{colsum}_u(G) < 1$, otherwise a direct rank test --- and
that the candidate $y^{\ast}$ is strictly positive;
the guard admits the branch on which the affine fixed
point is unique and positive, but does not by itself exclude further roots of
the nonlinear system; the multi-init admissibility checks recorded for the F1
column (three distinct starting points converging to the same root) are the
evidence that the reported root is the admitted one. Both the GAMMA and the DELTA solve
are subject to the same guard, and both converged to this root
(residuals $4.4\times10^{-16}$ to $2.0\times10^{-13}$
across the six fixed-wage cells,
\texttt{runs/matrix\_5x3-v10-\{GAMMA,DELTA\}-*/manifest.toml}). The
elasticity-independence is exact in the equations, not an $\epsilon$-limit
artefact: the DELTA solve at $\epsilon_\delta = 10^{-4}$ reproduces the
analytic $y^{\ast}$ with relative error $0.0$ (F2/F3 verification, recorded
in \texttt{docs/DOCS\_ASSESSMENT.md} section 4.2), and the v10 flow metrics
of the two rows agree to all printed decimals with residual-level wedge
differences $\sim 10^{-12}$.
\end{proof}

\begin{corollary}[Classification]\label{cor:class}
For the manuscript's framing of the matrix:

\begin{center}
\begin{tabular}{lll}
\toprule
Coincidence & Status & What is unidentified \\
\midrule
GAMMA $\equiv$ DELTA & definitional in form, & $(\theta,\epsilon,\sigma)$ \\
 & substantive in mechanism & in the fixed-wage gauge \\
 & (Lemma~\ref{lem:affine}) & \\
ALPHA $\equiv$ BETA & derivable & $\eta_s$ in the flexible-wage gauge \\
 & (Cor.~\ref{cor:realwage}) & \\
F2 $\equiv$ F3 (mobile rows) & derivable & the financing label; \\
 & (Section~\ref{sec:gauge}) & $F_{F3} = F_{F2} - B_{gov}$ \\
BF vs ALPHA & not a coincidence: & --- (the $\eta=0$ row is an \\
 & ADR-0020 option C gives the & identified equilibrium: own \\
 & frozen allocation a price & prices, wage vector \\
 & channel, so it is a distinct & $0.97$--$1.67$, own position) \\
 & economy & \\
sectoral corner $\equiv$ BF & exact, by construction & --- (the corner \emph{is} the \\
($\eta_{s,i} = 0$; Prop.~\ref{prop:nest}) & (Prop.~\ref{prop:nest}) & endpoint: the same equations) \\
\bottomrule
\end{tabular}
\end{center}

The first three entries are not numerical
accidents; they hold at machine precision in the executed generation, and
none is a mis-specification. They are properties of the \emph{experiment}: a
demand-only programme in a one-factor CRS economy cannot identify any
elasticity of the kernel (the mechanism is Remark~\ref{rem:samuelson}). The
fourth entry is different in kind: the BF/ALPHA coincidence held only in the
superseded $F = 0$ pin generation (ADR-0019 D2), where the frozen allocation
had no channel to prices; under ADR-0020 option C the $\eta = 0$ row carries
sectoral wages and is an identified equilibrium, so no BF quantity is left
unidentified and none must be withheld.
The fifth entry is an equivalence between two
\emph{closures} rather than between two rows of the executed matrix, and it is
exact by construction: the sectoral family's rigid corner reproduces the
$\eta = 0$ endpoint, which is why no separate rigid sectoral cell is reported.
\end{corollary}

\begin{remark}[Off the demand-only slice]\label{rem:supply}
The equivalences fail the moment technology moves. At $A \neq \one$,
Corollary~\ref{cor:realwage} gives $w/\Pi = 1/\Pi(\pi(A)) \neq 1$ in
general, and BETA's labour equation reduces to
$\alpha^{\!\top}y = \Lbar\,[1/\Pi(\pi(A))]^{\eta_s}$: employment is pinned
to a \emph{technology-determined constant} that now depends on $\eta_s$.
The supply-side pilot (+20\% in sector 1) measures exactly this separation
--- the failure of Proposition~\ref{prop:ab}'s hypothesis $A \equiv \one$:
$w = 1.004903$, with $L = 1.0024$ at $\eta_s = 0.5$ versus $L = 1.0098$ at
$\eta_s = 2$ against ALPHA's $L = 1.0000$ (\texttt{docs/ETAs.md}). The
identification of $\eta_s$ therefore requires a supply-side scenario ---
the premise of the supply arm.
\end{remark}

\subsection{The sectoral family and its rigid corner}\label{sec:sectoral}

The two equivalences above are coincidences \emph{within} the executed matrix.
This section records a third, of a different kind: an equivalence between two
\emph{closures}, one of which is the matrix's $\eta = 0$ endpoint and the other
a corner of a family that extends beyond the matrix. It is exact rather than
asymptotic, and it is what makes the labour door of Section~\ref{sec:scope}
concrete.

\paragraph{The family.} Replace the single aggregate supply equation
\eqref{eq:beta} by $N$ sectoral labour markets,
\begin{equation}
L^{cm}_i(p, y_i, w_i) \;=\; \bar{L}_i\,
   \Big[ \frac{w_i}{\Pi(p)} \Big]^{\eta_{s,i}},
\qquad i = 1, \dots, N, \qquad \eta_{s,i} \geq 0,
\label{eq:etas}
\end{equation}
solved jointly with zero profit \eqref{eq:zp}, clearing \eqref{eq:clearing},
the numeraire $\Pi(p) = 1$ and a free transfer $F$, in the unknowns
$X = (p, y, w(1{:}N), F) \in \mathbb{R}^{3N+1}$ (ADR-0022). The scalar closure
\eqref{eq:beta} is the member of this family with one common wage and
$\eta_{s,i} \equiv \eta_s$; the matrix's BETA row is therefore the scalar
\emph{special case} of a family, not a separate mechanism. Because the $N$
wages enter the $N$ zero-profit conditions \eqref{eq:zp} through the price
index \eqref{eq:ip}, the price block now depends on the $N$ employment levels
and Lemma~\ref{lem:price} no longer applies: demand reaches prices through the
wage structure, $p = \pi(A, w)$.

\paragraph{Two rigidities, one equation.} In both closures labour is immobile,
but for different reasons. At the $\eta = 0$ endpoint the allocation is frozen
\emph{by assumption}: $L_i = \bar{L}_i$ is a datum of the quantity block, and
the sectoral wage is whatever clears cost-minimizing demand against it. At
$\eta_{s,i} = 0$ each sector has a \emph{vertical} supply curve,
$L^s_i = \bar{L}_i$ for every $w_i$, so the market clears at the same
allocation with a demand-determined wage. A constraint imposed on a quantity
and an infinitely inelastic price response turn out to be the same row of
algebra --- which is the content of the proposition below.

\begin{proposition}[The rigid corner is the $\eta = 0$ endpoint]\label{prop:nest}
Set $\eta_{s,i} = 0$ for every $i$. Then the sectoral system \eqref{eq:etas}
and the $\eta = 0$ endpoint of ADR-0020 option C have the same solution set.
\end{proposition}

\begin{proof}
Write both systems in the unknowns $X = (p, y, w, F)$ as four row blocks:
\begin{align*}
\text{(P)} \quad & p_i = c_i(p, w_i; A), & & i = 1,\dots,N,
   & & \text{\eqref{eq:zp}} \\
\text{(Y)} \quad & y_i = int_i + c_i + (1-m_i)(g^G_i + inv_i + g_i) + X_i,
   & & i = 1,\dots,N, & & \text{\eqref{eq:clearing}} \\
\text{(W)} \quad & \log L^{cm}_i(p,y_i,w_i) - \log \bar{L}_i
   - \eta_{s,i} \log \frac{w_i}{\Pi(p)} = 0, & & i = 1,\dots,N, & & \\
\text{(E)} \quad & \Pi(p) = 1. & & & &
\end{align*}
The quantity block (Y) draws labour income from the wage bill
$\sum_i w_i L_i$ inside \eqref{eq:cons} and from the tax rate $\tau$; this is
the only place, besides (W), where the two systems differ. The $\eta = 0$
endpoint is the same four blocks with two changes: (W) drops the supply term
($\eta_{s,i} = 0$), and (Y) evaluates the wage bill at the \emph{frozen}
allocation, $L_i := \bar{L}_i$, instead of at the cost-minimizing demand
$L^{cm}_i$.

\emph{(i)} Let $X$ solve the sectoral system at $\eta_{s,i} \equiv 0$. Its (W)
gives $L^{cm}_i(X) = \bar{L}_i$ for every $i$, the supply term vanishing
identically. Substituting $L^{cm}_i = \bar{L}_i$ into its (Y) makes that block
equal the endpoint's (Y); (P) and (E) are identical in the two systems by
inspection. Hence $X$ solves the endpoint system.

\emph{(ii)} Let $X$ solve the endpoint system. Its (W) gives
$L^{cm}_i(X) = \bar{L}_i$, which satisfies the sectoral (W) at
$\eta_{s,i} \equiv 0$; substituting it into (Y) yields the sectoral (Y). Hence
$X$ solves the sectoral system. Both systems are square --- $3N+1$ unknowns,
$3N+1$ equations --- and their residual maps agree on the common solution set,
so the solution sets coincide.
\end{proof}

\begin{remark}[Why the family interpolates, and where it does not close]
Sweeping $\eta_{s,i}$ moves the adjustment from quantities to prices. At
$\eta_{s,i} = 0$ the allocation is pinned and the $N$ wages absorb the shock;
as $\eta_{s,i} \to \infty$ the supply curves become infinitely elastic, the
real wage returns to its technology-determined anchor and employment becomes
the absorbing margin --- GAMMA-like prices and GAMMA's absorbing margin. The
two ends of the family are therefore the matrix's own rigidity rows: BF
(allocation frozen, wages free) at the rigid corner, GAMMA (real wage pinned,
allocation free) at the elastic one. Two qualifications are exact and matter
for reading the table below. First, the elastic end is a \emph{limit}, not a
cell: the supply slope diverges and the direct solve eventually fails
($\eta_s = 10^6$; probe13). Second, the family keeps the transfer $F$ free,
whereas the fixed-wage system pins $F \equiv 0$; the family therefore
\emph{spans} the two rigidity rows without terminating in either.
\end{remark}

\paragraph{Measured.} Table~\ref{tab:sectoral} reports the executed sectoral cells
(\texttt{matrix\_5x3-v10-BETA-*}; ADR-0022). Four readings.

\emph{(a) Prices move, and with the demand composition.} Every sectoral cell
carries a deflator different from one and a $\max|p-1|$ that differs between F1
and F2/F3 --- the signature no single-wage closure produces (all twelve
$\eta = 1$ cells of the matrix read exactly $0.000000$).
\emph{(b) The elasticity splits the adjustment.} As $\eta_s$ rises from $0.25$
to $2$, the price response falls monotonically ($0.1529 \to 0.0372$ at F2)
while employment rises ($1.000864 \to 1.001914$): the family trades price
adjustment for quantity adjustment, and the welfare cost falls with it
($-1.69\% \to -1.39\%$). The F1 column, whose programme is financed by a
preference tilt rather than a tax, is the exception: its welfare index is
positive ($+0.42\%$ at $\eta_s = 0.5$).
\emph{(c) Incidence matters more than the level.} Making the \emph{programme}
sectors rigid ($\eta_{s,i} = 0$ there, $\eta_s = 0.5$ elsewhere) more than
doubles the price response ($0.2720$ against $0.1057$) and \emph{reverses} the
employment effect: employment falls by $0.34\%$ and the welfare index by
$2.69\%$, against $+0.13\%$ and $-1.57\%$ for the uniform cell. Rigidity on the
largest half by baseline employment instead gives $0.1234$ with a near-neutral
$+0.02\%$.
\emph{(d) Wage dispersion is large.} The mechanism moves relative wages:
$\max_i w_i / \min_i w_i$ reads $1.26$ at $\eta_s = 0.5$ and $1.70$ for the
rigid-programme variant, against $1.72$ at the $\eta = 0$ endpoint --- for a
programme worth $1.33\%$ of GDP.

\begin{table}[htbp]
\centering
\caption{The sectoral family in the \texttt{matrix\_5x3\_v10} generation
(\texttt{matrix\_5x3-v10-BETA-*}; deviations from the no-programme baseline;
F2 and F3 agree to all printed decimals in every cell).}
\label{tab:sectoral}
\begin{tabular}{lrrrr}
\toprule
Sectoral cell & $\max|p-1|$ & $\max|p-1|$ & Employment & Consumption \\
 & (F1) & (F2) & (F2) & (F2) \\
\midrule
uniform $\eta_s = 0.25$ & $0.133610$ & $0.152897$ & $1.000864$ & $-0.016888$ \\
uniform $\eta_s = 0.5$  & $0.095839$ & $0.105666$ & $1.001255$ & $-0.015712$ \\
uniform $\eta_s = 1$    & $0.061283$ & $0.065422$ & $1.001627$ & $-0.014669$ \\
uniform $\eta_s = 2$    & $0.035629$ & $0.037171$ & $1.001914$ & $-0.013912$ \\
\midrule
rigid programme, $\eta_s = 0.5$ & $0.215061$ & $0.272033$ & $0.996621$ & $-0.026914$ \\
rigid largest half, $\eta_s = 0.5$ & $0.109583$ & $0.123362$ & $1.000196$ & $-0.019053$ \\
\midrule
BF ($\eta_{s,i} = 0$) & $0.221442$ & $0.279009$ & $1.000000$ & $-0.019800$ \\
GAMMA (fixed wage) & $0.000000$ & $0.000000$ & $1.001260$ & $-0.015333$ \\
\bottomrule
\end{tabular}
\end{table}

\emph{The corner is the endpoint, numerically.} Solving the sectoral system at
$\eta_{s,i} = 0$ from the $\eta = 0$ solution reproduces $(p, y, w, F)$ exactly
--- $\max|\Delta| = 0.00 \times 10^0$ in all four blocks, and below $10^{-10}$
from a cold start (\texttt{tests/test\_sectoral\_labour.jl};
ADR-0022) --- as Proposition~\ref{prop:nest} requires. No separate
rigid sectoral cell is therefore reported: the corner \emph{is} the BF row of
Table~\ref{tab:matrix}.

\emph{The proof never uses $A = 1$.} Nothing in the four row blocks of the
proposition's proof refers to the technology shock, so the nesting is not a
demand-only property. Verified directly: under a $+20\%$ productivity shock in
sector 1 the canary again reads $0.00 \times 10^0$ in all four blocks
(\texttt{probe15\_sectoral\_supply.jl}; the figures are
recorded in \texttt{docs/log/2026-09.md} and \texttt{docs/ETAs.md}). The same probe shows why $\eta_{s,i}$
matters outside the demand-only slice: the rungs separate under the shock
(employment $1.000000$, $1.002187$, $1.008218$ and the real-wage spread
$1.215$, $1.105$, $1.043$ at $\eta_s = 0, 0.5, 2$), so the sectoral elasticity
is identifiable only with a supply-side scenario, while under a demand-only
programme it is a sensitivity parameter.


\subsection{What actually differs across the rows}\label{sec:gauge}

Lemma~\ref{lem:affine} shows that the demand-only matrix is a partition of
\emph{scalar closure regimes} over one and the same linear system.
Two scalar regimes occur across the $\eta = 1$
rows, and the $\eta = 0$ endpoint is now a third, genuinely nonlinear
regime (the sectoral-wage system):

\begin{center}
{\small
\begin{tabular}{lccc}
\toprule
Rows & $F$ & Employment & Household round in $G$ \\
\midrule
ALPHA, BETA ($\eta = 1$) & free (instrument) & pinned, $L = \Lbar$ & active \\
GAMMA, DELTA ($\eta = 1$) & $F \equiv 0$ & endogenous, $L = \alpha^{\!\top}y$ & active \\
BF ($\eta = 0$) & free (solved) & pinned via allocation & active (wage bill) \\
\bottomrule
\end{tabular}
}
\end{center}

\paragraph{The $F$-instrument.} In the mobile regime the labour pin selects
the transfer. Writing $y(F) = (I-G)^{-1}(b_0 + f\,F)$ with
$f = (1-s)\,\mathrm{diag}(1-m)\,\tilde\omega(ds)$, the pin
$\alpha^{\!\top}y = \Lbar$ gives
\begin{equation}
F^{\ast} = \frac{\Lbar - \alpha^{\!\top}(I-G)^{-1} b_0}
               {\alpha^{\!\top}(I-G)^{-1} f}.
\label{eq:finstrument}
\end{equation}
This instrument is why the mobile F2 and F3 columns are real-neutral: F2
charges the programme through the lump-sum tax ($E = \alpha^{\!\top}y -
\sum g^G - \sum g$), F3 through the external account ($E = \alpha^{\!\top}y
- \sum g^G + F$), and the two $b$-vectors differ by exactly
$f\,\sum g$. The pin \eqref{eq:finstrument} therefore yields
$F^{\ast}_{F3} = F^{\ast}_{F2} - \sum g$, identical real allocations, and
an identical net external position $F + B_{gov}$. Measured: $F_{F3} - F_{F2} = -0.01330991 = -G_0$ to eight
decimals in every mobile row (the $F$ values are printed
in \texttt{paper/tables/matrix\_5x3\_v6\_flows.md}, Table~2, which the v10
re-mint reproduces to all printed decimals; $G_0 = \sum_i g_i = 0.013310$,
ADR-0024).

\paragraph{Who pays.} The economically meaningful difference between the
two paradigms is which margin absorbs the programme. In the mobile rows the
labour pin forces the household to pay (consumption $-1.82\%$ under F2/F3;
the $F$ instrument crowds it out exactly), while real GDP is flat by
construction ($w\,\Lbar$ deflated, ADR-0018). In the fixed-wage rows
employment is the margin: F3 expands employment by $1.78\%$ and consumption
by $2.26\%$ (the Keynesian case), F2 leaves GDP nearly neutral
($+0.13\%$) with a welfare loss from the tax, and F1 shifts composition
slightly downward ($-0.10\%$ employment). The
$\eta = 0$ row is no longer degenerate: at the sectoral-wage endpoint the
welfare index reads $+0.096\%$ (F1) and $-1.980\%$ (F2/F3), the deflator
moves ($1.0063$ / $1.0072$), and the external account closes in every
column. Two caveats from the ADR-0020 audit now describe the
\emph{superseded} v5 pin generation: there the recorded BF/ALPHA
differences were carried by the pin, not by the frozen allocation
(replacing the pin by the labour equation reproduced the ALPHA cells to
$2.8\times10^{-17}$), and the $\eta = 0$ external position was not
identified (every pin value $F = c$ an exact root that moved the booked
position one-for-one). Option C resolves both: the account closes to
$\le 1.04\times10^{-11}$ at BF, $F$ is a solved equilibrium object,
financing neutrality extends to $\eta = 0$ ($F_{F3} = F_{F2} - B_{gov}$ to
$3.5\times10^{-14}$), and no BF entry need be withheld.

\paragraph{Admissibility.} All fifteen cells execute
(v6 and, re-minted, v10). In particular the F1 column is
admissible in the fixed-wage
rows: the saving rate and the import margins supply the
marginal leakage that keeps $\max \mathrm{colsum}(G) < 1$ (ADR-0017's
actual-matrix criterion reads $0.8314$ on full-71 and $0.8320$ under the F1
tilt); the endogenous tax rate $\tau = \sum g^G/\alpha^{\!\top}y$ is a
fixed-revenue level ($T = p\!\cdot\!g^G$) and enters $b$, not $G$ (the retired heuristic guard had blocked these
cells in the v1 generation).

\subsection{Scope conditions: what breaks the equivalences}\label{sec:scope}

Both propositions rest on Lemma~\ref{lem:price}: the price block contains
no demand variable. Any model change that lets demand into that block
restores identification and creates variety in the published matrix. The
doors are:

\begin{enumerate}
\item \textbf{Labour door --- sector-specific wages.} With $N$ sectoral
      labour markets $w_i = \phi_i(L_i)$ instead of one common $w$, the
      homogeneity argument of Lemma~\ref{lem:price} collapses: $L_i =
      \alpha_i y_i$ depends on demand, so $w_i$ and through zero profit
      $p_i$ become demand-sensitive. Zero profit is preserved (no rents),
      the new parameters are sectoral supply elasticities of the same
      family as $\eta_s$, and the variant repairs the BF row's
      common-wage defect (the \texttt{labor.BF} open
      gate of \texttt{registry/closures.toml}, resolved by ADR-0020 option C;
      \texttt{cbase2/review.md} section~1). Ranked candidate 1 in
      \texttt{docs/ETAs.md}. ADR-0020 option C is
      this door, adopted at the $\eta = 0$ endpoint: it is why the v6 BF
      row is a distinct economy rather than a normalisation of ALPHA.
      ADR-0021 (proposed) reuses the vector-wage apparatus in the
      fixed-wage row, but with an \emph{exogenous} pin --- that delivers
      design heterogeneity, not demand-sensitivity, because the pinned
      vector keeps the price block demand-free.
      ADR-0022 (accepted) generalises this door across
      the $\eta = 1$ rows: $N$ sectoral labour markets with supply elasticities
      $\eta_{s,i}$ (Section~\ref{sec:sectoral}), whose rigid corner is precisely
      the option C endpoint proved in Proposition~\ref{prop:nest}. It breaks the
      single-wage homogeneity on which Lemma~\ref{lem:price} rests, requires no
      income closure (no rents are created), and its scalar member is the
      existing BETA row.
\item \textbf{Markup door --- demand-sensitive markups.} A pricing rule
      $p_i = \mu_i(y_i)\, c_i(p,w)$, e.g.\ $\mu_i = 1 + \kappa\,(y_i/y_{i0}
      - 1)$, breaks Lemma~\ref{lem:price} directly. The real cost is not
      code but accounting: markup income is profit, so a
      profit-income-and-spending closure is required or the external-account
      identity behind the canary breaks. Nested at $\kappa = 0$. Candidate 2
      in \texttt{docs/ETAs.md}.
\item \textbf{Capacity door --- a fixed factor, or its reduced form.} A
      sector-fixed factor with a rental schedule makes marginal cost rise
      in $y_i$ (rents require a rental-income closure --- the largest
      surgery, candidate 3 in \texttt{docs/ETAs.md}). The reduced-form
      variant avoids all income-closure problems: an endogenous utilization
      externality $A^{eff}_i = A_i\,(y_i/\lambda_i)^{-\delta}$ raises unit
      cost with demand while keeping zero profit (the extra cost is paid to
      no one), so a one-line change in the cost hook \eqref{eq:cost}
      suffices. With the opposite sign ($\delta < 0$) it is the
      Kaldor--Verdoorn case: increasing returns make prices \emph{fall}
      with demand --- heterodox-friendly variety in the other direction.
      (New candidate; see Section~\ref{app:eqnumevidence}.)
\item \textbf{External door --- endogenous import prices.} An
      upward-sloping foreign supply curve
      $p^M_i = p^M_{i0}\,(M_i/M_{i0})^{\zeta}$, $\zeta \geq 0$, enters the
      intermediate index \eqref{eq:ip} and the CPI through the margins, so
      $\pi$ becomes a function of import volumes and hence of demand. The
      surgery is moderate and it exploits precisely the external-account
      machinery the paper has just built (ADR-0019): the programme's import
      content would bid up import prices --- a structuralist
      absorption/term-of-trade channel that fits the open-economy framing.
      (New candidate.)
\item \textbf{Technology door --- supply shocks.} $A \neq 1$ breaks
      Corollary~\ref{cor:realwage} and separates BETA from ALPHA
      (Remark~\ref{rem:supply}); this is the supply arm already designed in
      \texttt{docs/ETAs.md}, not a repair of the demand-only matrix.
\end{enumerate}

\begin{corollary}[No-go for nominal wage rules]\label{cor:nogo}
Within the kernel's one-factor CRS price block, no rule that writes the
\emph{nominal} wage as a function of endogenous aggregates --- an aggregate
wage curve $w = \bar{w}\,(\alpha^{\!\top}y/\Lbar)^{\varphi}$, nominal
indexation, or a re-anchored numeraire --- can create demand-sensitive
\emph{real} outcomes. By Lemma~\ref{lem:price}, $p = w\,\pi(A)$ and
$\Pi(p) = w\,\Pi(\pi)$ at every solution, so the real wage
$w/\Pi = 1/\Pi(\pi(A))$ is invariant to whatever sets $w$; the rule merely
re-labels the nominal scale. With the CPI numeraire or the $w = 1$ pin such
a rule is either vacuous or over-determines the system (in the fixed gauge
it forces $\alpha^{\!\top}y = \Lbar$, i.e.\ silently reverts GAMMA to
ALPHA). A rule on the \emph{real} wage, $w/\Pi = \phi(L)$, is BETA's
equation \eqref{eq:beta} rearranged and is already in the model --- where
it is pinned by technology under demand-only shocks. The numeraire lever is
likewise dead (DE-0011).
A single common wage curve is also what a
\emph{single-wage} price block can express; the sectoral-wage door above
escapes the corollary precisely by breaking the single-wage homogeneity on
which it rests.
\end{corollary}

Quantity-side levers (endogenous saving rates, endogenous export demand)
change $G$ or $b$ but not the price block; they add variety between
financing columns, not between closures, and leave Lemma~\ref{lem:price}
intact.

\subsection{Numerical evidence}\label{app:eqnumevidence}

The fifteen matrix cells were re-executed as the
\texttt{matrix\_5x3\_v9} generation on the extended kernel that carries the
sectoral family (ADR-0022), and are cited here from the
re-minted \texttt{matrix\_5x3\_v10} generation (the
C1-corrected measurement, 2026-09-20). v6 and v9 are bit-identical to each
other ($\max|\Delta| = 0$ on \texttt{gdp\_rel}, \texttt{consumption\_rel},
employment, the deflator, $\max|p-1|$, the external transfer, the wage extremes
and the identity residual); v10 reproduces them to solver precision (max
$9.34\times10^{-9}$ at the stiff BF-F1 \texttt{wage\_max} diagnostic,
$1.63\times10^{-11}$ in the worst ALPHA/BETA/GAMMA/DELTA row,
$4.55\times10^{-15}$ at the exact-baseline cells; no printed figure is
affected, \texttt{paper/tables/matrix\_5x3\_v10\_flows.md}, Table~2), so the
equivalences and the table below are invariant to the extension. The sectoral
cells of the same generation are reported separately in
Table~\ref{tab:sectoral}.

Table~\ref{tab:matrix} extracts the rows of the executed
\texttt{matrix\_5x3\_v10} generation (the fifteen matrix
cells reproduce their v6/v9 manifests to all printed decimals) (income-side real GDP and the
household-consumption welfare index, ADR-0018; employment). The paired rows
(ALPHA/BETA, GAMMA/DELTA) are identical to all printed decimals in every
reported metric, including the full flow decomposition (saving, tax,
investment plus exports, imports, taxes on intermediates; see
\texttt{paper/tables/matrix\_5x3\_v6\_flows.md}, whose matrix rows the v9 and v10
generations reproduce to all printed decimals). The
three BF rows are included for contrast: the $\eta = 0$ sectoral-wage
endpoint is a distinct economy, not a normalisation of ALPHA. The GDP
deflator reads $1.000000$ in the twelve $\eta = 1$
cells --- the numerical face of Lemma~\ref{lem:price} --- while the three BF
$\eta = 0$ cells carry a deflator of $1.006326$ (F1) and $1.007181$
(F2/F3), the price response the sectoral wages create.

\begin{table}[htbp]
\centering
\caption{The equivalence rows in the \texttt{matrix\_5x3\_v10} generation
(run ids \texttt{matrix\_5x3-v10-*}; deviations from the no-programme
baseline).}
\label{tab:matrix}
\begin{tabular}{llrrr}
\toprule
Row & Financing & Real GDP & Consumption (welfare) & Employment $L$ \\
\midrule
ALPHA & F1 & $+0.000000$ & $+0.001432$ & $1.000000$ \\
BETA  & F1 & $+0.000000$ & $+0.001432$ & $1.000000$ \\
ALPHA & F2 & $+0.000000$ & $-0.018226$ & $1.000000$ \\
BETA  & F2 & $+0.000000$ & $-0.018226$ & $1.000000$ \\
ALPHA & F3 & $+0.000000$ & $-0.018226$ & $1.000000$ \\
BETA  & F3 & $+0.000000$ & $-0.018226$ & $1.000000$ \\
\midrule
GAMMA & F1 & $-0.000973$ & $-0.000806$ & $0.999027$ \\
DELTA & F1 & $-0.000973$ & $-0.000806$ & $0.999027$ \\
GAMMA & F2 & $+0.001260$ & $-0.015333$ & $1.001260$ \\
DELTA & F2 & $+0.001260$ & $-0.015333$ & $1.001260$ \\
GAMMA & F3 & $+0.017796$ & $+0.022644$ & $1.017796$ \\
DELTA & F3 & $+0.017796$ & $+0.022644$ & $1.017796$ \\
\midrule
BF & F1 & $-0.000110$ & $+0.000963$ & $1.000000$ \\
BF & F2 & $-0.000162$ & $-0.019800$ & $1.000000$ \\
BF & F3 & $-0.000162$ & $-0.019800$ & $1.000000$ \\
\bottomrule
\end{tabular}
\end{table}

\paragraph{Implications for the manuscript.} (i) The two-paradigm reading
is safe to state as a theorem-backed claim: the flexible-wage rows form one
equivalence class (composition at pinned employment), the fixed-wage rows
another (employment at pinned real wage), and the difference between the
classes is the $F$-instrument --- a derivation, not an observation.
(ii) The null results for $\eta_s$ and for $(\theta,\epsilon,\sigma)$ under
demand-only shocks should be reported as identification statements
(Corollary~\ref{cor:class}), which is precisely what reviewer 2's critique
needs. (iii) The variety program for the demand-only matrix must route
through one of the five doors of Section~\ref{sec:scope}; within it, the
capacity door's utilization/Verdoorn variant and the external door are the
cheap additions beyond the three candidates already ranked in
\texttt{docs/ETAs.md}, while Corollary~\ref{cor:nogo} rules out the entire
nominal-wage-rule family ex ante. The labour door is
no longer hypothetical: ADR-0020 option C delivers it at the $\eta = 0$
endpoint, and the general sectoral-labour-market form extends it across the
$\eta = 1$ rows.
(iv) The labour door is now an implemented mechanism
rather than a proposal: the sectoral family of Section~\ref{sec:sectoral} is
promoted, tested and executed (ADR-0022), it delivers demand-sensitive prices
inside a flexible-wage closure, and it leaves the executed matrix
unchanged to solver precision (no printed figure
affected). The manuscript can therefore state three exact equivalences and
present the family as the answer to the variety question --- with two caveats
that must travel with it: under demand-only shocks $\eta_{s,i}$ is a
\emph{sensitivity} parameter, not an estimate, so the level of the price
response is identified only by the supply arm (Remark~\ref{rem:supply}); and
the rule that makes a subset of sectors rigid is a modelling choice, to be
reported as a band rather than a point.

% =====================================================================
% Appendix B --- Kernel, calibration, A-bill accounting
% ---------------------------------------------------------------------
% Job: Destatis rows 73-75, the external-account canary identity; plus the
%      reproducibility subsection (run-id citation discipline, the v9/v10
%      split, the preregistration trail with 41 failed / 3 provisional runs
%      left visible --- a credibility asset for this submission).
% Mindmap: N05.
% Status: STUB --- narrative outline v3 (2026-09-20).
% =====================================================================
\section{Kernel, calibration, and A-bill accounting}\label{app:kernel}

% TODO: the common kernel and its invariants; the A-bill intermediate-bill
% accounting (Destatis rows 73/74/75: domestic bill A, imported intermediates
% M_int, product taxes T_int; identity sum A + sum M_int + sum T_int =
% sum lambda - 1 to 2e-16); the external-account canary identity
% S + T_int + M - (I + X) = F + B_gov.

\subsection*{Reproducibility}
% TODO: every table cites run_ids; the v9/v10 citation split (v10 for the
% matrix and the 18 sectoral cells, v9 for everything else); the
% preregistration trail (373 executed, 41 failed, 3 provisional, left
% visible per ADR-0004).
% =====================================================================
% Appendix C --- Matrix variants and gamma/beta representation robustness
% ---------------------------------------------------------------------
% Job: where the presentation-decision variants live; all tables cite run
%      ids. Tally: fifteen matrix cells + twelve ladder cells + six
%      two-group cells = 33.
% Mindmap: N11, N12.
% Status: STUB --- narrative outline v3 (2026-09-20).
% =====================================================================
\section{Matrix variants and gamma/beta representation robustness}\label{app:variants}

% TABLE (linked, revision/tables/matrix_5x3_v10_flows.md): the full 15-cell
% matrix + the 18 sectoral cells, all run ids.
% TODO: the full 15-cell matrix (citing matrix_5x3_v10 run ids); the
% gamma/beta variant tables; the ADR-0021 wage-structure axis as a
% representation-robustness row; the 33-cell tally
% (15 matrix + 12 ladder + 6 two-group).
% =====================================================================
% Appendix D --- Vector construction and data provenance
% ---------------------------------------------------------------------
% Job: impulses.csv, the ADR-0024 price-basis verification.
% Mindmap: N08.
% Status: STUB --- narrative outline v3 (2026-09-20).
% =====================================================================
\section{Vector construction and data provenance}\label{app:provenance}

% TODO: construction of the green-investment vector from impulses.csv
% (27 rows x 75 columns); the 2024 row (80.76 bn, current/2023 prices) and
% the horizon mean; the deflator 1.46 (construction-cost indices); the
% ADR-0024 verification (58.89 / 1.46 = 40.338 = 40.337 bn); the incidence
% psi (raw 2024 current-price vector) and the open composition item.
% =====================================================================
% Appendix E --- Closures and the neo-Keynesian mapping
% ---------------------------------------------------------------------
% Ported from paper/to_evaluate/closures_and_demand_shocks.md (2026-09-18).
% This appendix is NEW relative to the outline's A-D plan; it carries the
% CGE-vs-IO cleavage, the FIDELIO mapping and the closure taxonomy that
% Section 3 introduces.
%
% *** REVIEWER FLAG (Stage-3 must-not-appear list) ***
% The closing subsection's dominance wording ("the wage regime is the
% first-order determinant of the aggregate multiplier"; "intersectoral
% mobility eta produces second-order reallocation effects") sits close to
% the retired eta-dominance claims. Check it against the must-not-appear
% list before this text enters the manuscript. The descriptive mapping
% (Sections 1-4) is uncontroversial; the summary's two claims are not.
% =====================================================================
\section{Closures and the neo-Keynesian mapping}\label{app:mapping}

\noindent This appendix maps the CGE/IO cleavage and the neo-Keynesian
alternative --- the EU JRC FIDELIO model --- onto the closure matrix. It
supports Section~\ref{sec:two-questions}.

\subsection{The core theoretical cleavage: demand shocks in CGE versus IO models}\label{app:cleavage}

A central question in applied input--output modeling --- especially for
evaluating green public investment programmes such as the EU Recovery and
Resilience Facility or large-scale building renovation --- is whether an
exogenous demand stimulus:

\begin{enumerate}
\item \textbf{expands real aggregate output} by mobilising unutilised
  resources (the Keynesian / Leontief perspective), or
\item \textbf{merely reallocates scarce resources}, causing factor prices to
  rise and crowding out private activity with virtually zero net change in
  aggregate real GDP (the canonical neoclassical CGE perspective).
\end{enumerate}

\paragraph{Why standard CGE models produce minimal aggregate effects from demand shocks.}
In textbook Walrasian CGE models: (i) \emph{full factor utilisation} ---
aggregate primary factor endowments $(\bar L, \bar K)$ are exogenously given
and fully employed; (ii) \emph{instantaneous market clearing} --- prices and
wages are perfectly flexible and clear all commodity and factor markets
instantaneously; (iii) \emph{Say's Law operates} --- output is
supply-determined by the production possibility frontier
$Y = F(K, L, \mathbf{X})$; (iv) \emph{crowding out} --- an exogenous increase
in public investment raises demand in targeted sectors, but because total
labour and capital cannot expand, the extra factor demand bids up wages and
capital rentals, and the relative price increase squeezes the margins of
non-shocked sectors, crowding out private consumption, private investment or
net exports. The aggregate real GDP response is negligible
($\Delta Y \approx 0$).

\subsection{The macroeconometric alternative: the EU JRC FIDELIO model}\label{app:fidelio}

The \textbf{FIDELIO} model (\emph{Fully Interregional Dynamic Econometric
Long-term Input--Output}; JRC / DG GROW; Rocchi et al., 2019, 2025) was built
specifically to overcome this limitation. Rather than adopting a Walrasian
general-equilibrium core, FIDELIO is a \emph{dynamic macroeconometric,
multi-sectoral Neo-Keynesian} model grounded in the official Eurostat/JRC
FIGARO inter-country input--output tables (64 NACE industries, 46
countries/regions). It generates positive, non-trivial real output and
employment multipliers from demand shocks through five structural features:

\begin{enumerate}
\item \textbf{Involuntary unemployment and the wage-setting curve.} Wages do
  not clear the labour market. The baseline economy features involuntary
  unemployment ($u > 0$). Nominal wages adjust sluggishly along an
  econometrically estimated \emph{wage curve} (Blanchflower--Oswald
  tradition):
  \[
  \Delta \ln w_{i,r,t} = f(\text{inflation}_{t-1},\ \text{productivity}_{t},\ u_{r,t-1}).
  \]
  Because labour supply is not vertical at baseline, firms can hire idle
  workers to expand production without hitting an immediate supply wall.
\item \textbf{Mark-up pricing and sluggish cost pass-through.} Product prices
  are set via markups over unit costs (derived from flexible Translog cost
  functions over capital, labour, energy and intermediate inputs) rather than
  jumping to Walrasian clearing levels. Output in the short-to-medium run is
  \emph{demand-determined}.
\item \textbf{Keynesian multiplier loops (the ECM consumption block.)}
  Household consumption follows a dynamic error-correction model tied to
  current real disposable income rather than to an intertemporal Euler
  equation:
  \[
  \Delta G \Rightarrow \uparrow Y \Rightarrow \uparrow L
  \Rightarrow \uparrow (w\cdot L) \Rightarrow \uparrow C^{\text{household}},
  \]
  inducing a secondary Keynesian expenditure loop (a Type-II multiplier),
  yielding aggregate output multipliers between 1.2 and 1.8 in JRC policy
  simulations.
\item \textbf{Dynamic capital accumulation (endogenous capacity).} Capital is
  not a fixed static endowment; investment responds dynamically through an
  accelerator / user-cost formulation, so sustained demand increases capacity
  utilisation and expected returns, prompting firms to invest and expand
  capital stocks over time.
\item \textbf{Medium-to-long-run re-equilibration.} As unemployment drops and
  capacity tightens, the wage curve drives up wages and unit costs; the
  resulting rise in output prices erodes export competitiveness, generating
  gradual crowding out toward a medium-run NAIRU --- unless the stimulus was
  in public capital or R\&D, in which case it permanently shifts potential
  output outward.
\end{enumerate}

\subsection{Mapping into the \texttt{BeyondHulten} $5\times3$ closure matrix}\label{app:matrix-map}

The platform organises macroeconomic closures across two dimensions (five
labour closures $\times$ three financing closures).

\paragraph{The five labour closures.}
\begin{enumerate}
\item \textbf{ALPHA} (full-employment mobile CGE benchmark):
  $\sum_i L_i = \bar L$, single flexible wage $w$, cost-minimising
  intersectoral allocation ($\eta = 1$).
\item \textbf{BF} (immobile-labour endpoint selector): $\eta = 0$ freezes
  sectoral labour quantities to baseline shares $L_i = \bar L_i$; with the
  sectoral-wage endpoint (ADR-0020) the wage vector is free and the frozen
  allocation carries a price channel.
\item \textbf{BETA} (elastic labour supply on the real wage):
  $\sum_i L_i = \bar L\,[(w/P)/(w_0/P_0)]^{\eta_s}$, deflated by the CPI.
\item \textbf{GAMMA} (fixed real wage, uncapped extensive margin):
  $w/P = \bar w = 1$, employment $L = \sum_i L_i$ completely
  demand-determined.
\item \textbf{DELTA} (Leontief IO corner): GAMMA ($w/P = 1$) combined with
  the zero-substitution limit of the CES production core
  ($\theta, \epsilon, \sigma \to 0^+$), reproducing the fixed-price Leontief
  multiplier.
\end{enumerate}
(Complementarity extension \textbf{ZETA}: $0 \le \bar L - L \perp w/P -
\bar\omega \ge 0$, regime-switching between GAMMA slack and ALPHA full
employment; recorded as an idea, not a closure.)

\paragraph{The three financing closures.}
\begin{enumerate}
\item \textbf{F1} (preference reallocation / demand tilt): budget-neutral
  composition shift within fixed household expenditure
  ($\sum_i p_i c_i = E_h$); no net fiscal injection.
\item \textbf{F2} (tax-financed public investment): balanced government
  budget ($\sum_i p_i g_i = T(p)$); autonomous public investment financed by
  an explicit domestic tax counterparty.
\item \textbf{F3} (external / debt-financed programme): autonomous public
  spending financed externally or via debt ($\sum_i p_i g_i = F$), with
  imports and external balances absorbing the deficit.
\end{enumerate}

\paragraph{Matrix location of models.}
\begin{table}[htbp]
\centering
\small
\begin{tabular}{l l l l}
\toprule
Labour row & F1 (demand tilt) & F2 (tax-financed) & F3 (debt / external) \\
\midrule
ALPHA (full $L$)   & sectoral reallocation & full crowding out & crowding out ($w$ bids up) \\
BF (immobile)      & bottlenecks / wedges  & contractionary wedge & bottleneck limits \\
BETA (elastic)     & moderate shift        & mild expansion & expansion via $w/P$ \\
GAMMA (fixed $w$)  & demand-driven shift   & balanced-budget multiplier & Keynesian multiplier \\
DELTA (IO/Leontief)& IO composition shift  & Leontief balanced multiplier & pure IO multiplier \\
\bottomrule
\end{tabular}
\caption{Where the standard CGE and the FIDELIO model sit in the matrix.}
\label{tab:matrix-map}
\end{table}

\begin{itemize}
\item \textbf{Where standard CGE sits:} ALPHA-F1 or ALPHA-F2 --- fixed
  aggregate labour $\bar L$ and flexible wages force crowding out; real output
  effects are close to zero.
\item \textbf{Where FIDELIO sits:} in its static / basic IO mode (modules
  1--2) at DELTA-F3 (a pure Leontief input--output multiplier driven by
  external funding); in its short-run dynamic mode at GAMMA-F3 (fixed real
  wage, unconstrained extensive margin, autonomous injection), activating the
  full Keynesian respending multiplier; and, in the medium-to-long run, moving
  dynamically from GAMMA-F3 toward BETA-F3 as the wage curve drives real wage
  growth.
\end{itemize}

\subsection{How neo-Keynesian assumptions map into our closures}\label{app:neokeynesian}

In this platform, neo-Keynesian assumptions are not auxiliary parameters ---
they define the mathematical structure of the closure.

\paragraph{A. Wage rigidity (extensive margin versus wage inflation).}
In accordance with ADR-0014 and DE-0004, wage rigidity here is defined on the
\emph{real wage} ($w/P = \bar w$), deflated by the CPI numeraire ($P = 1$).
In GAMMA, fixing $w/P = 1$ removes the wage as an equilibrating price: the
labour-market equation changes from clearing $L = \bar L$ to an unconstrained
demand equation $L = \sum_i L_i(y_i, \mathbf{p})$, and the extensive margin of
employment absorbs the demand shock. Asymmetric rigidity (ZETA) follows the
spirit of Baqaee--Farhi (2022): downward wage rigidity creates a regime
switch --- demand deficits cause unemployment (GAMMA regime), while demand
surges that exhaust slack trigger wage inflation (ALPHA regime).

\paragraph{B. Price rigidity and fixed markups.} In standard New-Keynesian
models, Calvo or Rotemberg pricing prevents prices from adjusting
instantaneously to demand shocks. Here, under GAMMA, wages are pegged
($w = 1$) and the shock is demand-driven, so marginal costs remain roughly
constant and product prices exhibit high inertia ($\mathbf{p} \approx 1$). In
DELTA, the zero-substitution limit ($\theta, \epsilon, \sigma \to 0^+$)
prevents firms from adjusting technical factor proportions, so cost
pass-through becomes purely linear and additive, reproducing the classic
fixed-price Leontief/Robinson (2006) multiplier.

\paragraph{C. Non-Ricardian demand and disposable-income loops.} Standard CGE
models enforce intertemporal optimisation in which households anticipate
future tax liabilities, neutralising debt-financed spending. Here the
consumption block routes household income $E_h = (1-\tau) w \sum_i L_i +
\text{capital income}$; under F3 the absence of an immediate tax hike allows
the newly generated wage income $w\cdot\Delta L$ to fuel private consumption
directly, generating an endogenous Keynesian multiplier.

\subsection{Summary and methodological takeaway}\label{app:summary}
% *** Flagged above: verify against the must-not-appear list. ***

\begin{enumerate}
\item \textbf{The binary wage regime dominates.} The literature often debates
  the elasticity of factor substitution or intersectoral mobility ($\eta$).
  As shown in \texttt{definitive\_guide.md} and our runs, intersectoral
  mobility produces second-order reallocation effects, whereas the wage regime
  (flexible ALPHA versus sticky GAMMA) is the first-order determinant of the
  aggregate multiplier.
\item \textbf{Reconciling applied discrepancies.} When policy studies using
  macroeconometric models (FIDELIO, E3ME) report output multipliers around
  $1.5$ for green investments while standard CGE models report zero or
  negative multipliers, the divergence is not an empirical dispute over green
  technology. It is a direct mathematical consequence of evaluating the shock
  at GAMMA/DELTA-F3 (slack resources, wage stickiness, debt financing) versus
  ALPHA-F2 (full employment, market clearing, domestic tax financing).
\end{enumerate}
\printbibliography[title={Appendix references}]

\end{refsection}

\ifendfloats
  \renewcommand{\thetable}{A\arabic{table}}
  \renewcommand*{\theHtable}{\thetable}
  \renewcommand{\thefigure}{A\arabic{figure}}
  \renewcommand*{\theHfigure}{\thefigure}
  \renewcommand{\thepostfigure}{A\arabic{postfigure}}
  \renewcommand{\theposttable}{A\arabic{posttable}}
  \renewcommand*{\theHposttable}{appendix.posttable.\arabic{posttable}}
  \processdelayedfloats
\fi

\end{document}